\documentclass[JEP,XML,SOM,Unicode,published]{cedram}
\datereceived{2020-10-20}
\dateaccepted{2021-04-30}
\dateepreuves{2021-05-10}

\usepackage{mathrsfs}
\let\scr\mathscr

\newcommand\mto{\mathchoice{\longmapsto}{\mapsto}{\mapsto}{\mapsto}}
\newcommand\hto{\mathrel{\lhook\joinrel\to}}
\newcommand{\To}[1]{\mathchoice{\xrightarrow{\ts\kern4pt#1\kern3pt}}{\stackrel{#1}{\longrightarrow}}{}{}}
\let\ra\rightarrow

\let\vref\ref
\let\ts\relax
\newcommand{\sbullet}{{\scriptscriptstyle\bullet}}
\newcommand{\ab}{\mathrm{ab}}
\newcommand{\an}{\mathrm{an}}

\newenvironment{enumeraten}
{\begin{enumerate}\setcounter{enumi}{\value{equation}}\itemindent0pt}
{\setcounter{equation}{\value{enumi}}\end{enumerate}}

\usepackage{newunicodechar}
\usepackage{mathtools}
\usepackage{tikz}
\usepackage{tikz-cd}
\usepackage{nicematrix}

\tikzcdset{arrow style=Latin Modern}
\newcommand{\dradpl}{{\hspace*{-2.5mm}\begin{tikzcd}[ampersand replacement=\&,column sep = scriptsize]\mbox{}\ar[r,dashed]\&\mbox{}\end{tikzcd}\hspace*{-2.5mm}}}
\newcommand{\dratst}{{\hspace*{-2.5mm}\begin{tikzcd}[ampersand replacement=\&,column sep = small]\mbox{}\ar[r,dashed]\&\mbox{}\end{tikzcd}\hspace*{-2mm}}}
\let\drascst\dratst
\newcommand{\dra}{\mathchoice{\dradpl}{\dratst}{\drascst}{\drascst}}
\let\dasharrow\dra

\DeclareFontFamily{U}{mathx}{\hyphenchar\font45}
\DeclareFontShape{U}{mathx}{m}{n}{
<5> <6> <7> <8> <9> <10>
<10.95> <12> <14.4> <17.28> <20.74> <24.88>
mathx10
}{}
\DeclareSymbolFont{mathx}{U}{mathx}{m}{n}
\DeclareFontSubstitution{U}{mathx}{m}{n}
\DeclareMathAccent{\wcheck}{0}{mathx}{"71}

\theoremstyle{plain}
\newtheorem{thm}{Theorem}[section]
\newtheorem{claim}[thm]{Claim}
\newtheorem{consequence}[thm]{Consequence}
\newtheorem{fact}[thm]{Fact}
\newtheorem{lem}[thm]{Lemma}
\newtheorem{prop}[thm]{Proposition}

\theoremstyle{definition}
\newtheorem{defn}[thm]{Definition}
\newtheorem{asswlog}[thm]{Assumption w.l.o.g.}
\newtheorem{construction}[thm]{Construction}
\newtheorem{notation}[thm]{Notation}
\newtheorem{rem}[thm]{Remark}

\numberwithin{equation}{thm}
\numberwithin{enumi}{thm}

\DeclareMathOperator{\Aut}{Aut}
\DeclareMathOperator{\codim}{codim}
\DeclareMathOperator{\img}{img}
\DeclareMathOperator{\Pic}{Pic}
\DeclareMathOperator{\rank}{rank}
\newcommand{\reg}{\mathrm{reg}}
\DeclareMathOperator{\sEnd}{\sE\negthinspace \mathit{nd}}
\DeclareMathOperator{\sing}{sing}
\DeclareMathOperator{\Sym}{Sym}
\DeclareMathOperator{\supp}{supp}
\DeclareMathOperator{\tor}{tor}

\DeclareMathOperator{\alb}{alb}
\DeclareMathOperator{\Alb}{Alb}
\DeclareMathOperator{\diff}{d}
\DeclareMathOperator{\Div}{Div}
\DeclareMathOperator{\Gal}{Gal}
\DeclareMathOperator{\GL}{GL}
\DeclareMathOperator{\iit}{iit}
\DeclareMathOperator{\Iit}{Iit}
\DeclareMathOperator{\Lie}{Lie}
\DeclareMathOperator{\PGL}{\bP\negthinspace\GL}
\DeclareMathOperator{\Rad}{Rad}
\DeclareMathOperator{\sh}{sha}
\DeclareMathOperator{\Sh}{Sha}
\DeclareMathOperator{\SL}{SL}

\newcommand{\sA}{\scr{A}}
\newcommand{\sE}{\scr{E}}
\newcommand{\sF}{\scr{F}}
\newcommand{\sL}{\scr{L}}
\newcommand{\sN}{\scr{N}}
\newcommand{\sO}{\scr{O}}
\newcommand{\sS}{\scr{S}}
\newcommand{\sT}{\scr{T}}
\newcommand{\sU}{\scr{U}}

\newcommand{\bC}{\mathbb{C}}
\newcommand{\bN}{\mathbb{N}}
\newcommand{\bP}{\mathbb{P}}
\newcommand{\bQ}{\mathbb{Q}}
\newcommand{\bZ}{\mathbb{Z}}

\let\wtilde\widetilde
\let\what\widehat

\hyphenation{com-po-nents}
\hyphenation{pos-i-tive}
\hyphenation{Theo-rem}
\hyphenation{Vojta}

\newcommand{\factor}[2]{\left. \raise 2pt\hbox{$#1$} \right/\hskip -2pt\raise -2pt\hbox{$#2$}}

\datepublished{2021-05-20}
\begin{document}
\frontmatter
\title{Projectively flat klt varieties}

\author[\initial{D.} \lastname{Greb}]{\firstname{Daniel} \lastname{Greb}}
\address{Essener Seminar für Algebraische Geometrie und Arithmetik, Fakultät für Mathematik, Universität Duisburg--Essen\\
45117 Essen, Germany}
\email{daniel.greb@uni-due.de}
\urladdr{https://www.esaga.uni-due.de/daniel.greb}

\author[\initial{S.} \lastname{Kebekus}]{\firstname{Stefan} \lastname{Kebekus}}
\address{Mathematisches Institut, Albert-Ludwigs-Universität Freiburg\\ Ernst-Zermelo-Straße 1, 79104 Freiburg im Breisgau, Germany}
\address{\& Freiburg Institute for Advanced Studies (FRIAS)\\
Freiburg im Breisgau, Germany}
\email{stefan.kebekus@math.uni-freiburg.de}
\urladdr{https://cplx.vm.uni-freiburg.de}

\author[\initial{T.} \lastname{Peternell}]{\firstname{Thomas} \lastname{Peternell}}
\address{Mathematisches Institut, Universität Bayreuth\\
95440~Bayreuth, Germany}
\email{thomas.peternell@uni-bayreuth.de}
\urladdr{https://www.komplexe-analysis.uni-bayreuth.de}

\thanks{Daniel Greb is partially supported by the DFG-Research Training Group
``Symmetries and classifying spaces: analytic, arithmetic and derived''.
Stefan Kebekus gratefully acknowledges partial support through a fellowship of
the Freiburg Institute of Advanced Studies (FRIAS). Thomas Peternell is
partially supported by a DFG grant ``Zur Positivität in der Komplexen
Geometrie'' and gratefully also acknowledges the support by FRIAS}

\begin{abstract}
In the context of uniformisation problems, we study projective varieties with klt singularities whose cotangent sheaf admits a projectively flat structure over the smooth locus. Generalising work of Jahnke-Radloff, we show that torus quotients are the only klt varieties with semistable cotangent sheaf and extremal Chern classes. An analogous result for varieties with nef normalised cotangent sheaves follows.
\end{abstract}

\subjclass{32Q30, 32Q26, 14E20, 14E30, 53B10}

\keywords{Bogomolov-Gieseker inequality, Abelian variety, klt singularities, Miyaoka-Yau inequality, stability, projective flatness, uniformisation}

\alttitle{Variétés klt projectivement plates}

\altkeywords{Inégalité de Bogomolov-Gieseker, variété abélienne, singularités klt, inégalité de Miyaoka-Yau, stabilité, platitude projective, uniformisation}

\begin{altabstract}
Dans le cadre des problèmes d'uniformisation, nous étudions les variétés projectives avec singularités klt dont le faisceau cotangent admet une structure projective plate sur le lieu lisse. En généralisant le travail de Jahnke-Radloff, nous montrons que les quotients des tores sont les seules variétés klt avec un faisceau cotangent semi-stable et des classes de Chern extrémales. Un résultat analogue pour les variétés avec un faisceau cotangent normalisé nef s'ensuit.
\end{altabstract}

\maketitle
\tableofcontents
\mainmatter

\section{Introduction}

\subsection{Projective manifolds with projectively flat cotangent bundle}

Let $\sE$ be a locally free sheaf of rank $r \geq 1$ on a complex, projective
manifold $X$ of dimension $n \geq 2$. If $\sE$ is semistable with respect to some
ample divisor $H \in \Div(X)$, then the Bogomolov-Gieseker Inequality holds:
\begin{equation}\label{eq:x1}
\frac{r-1}{2r}\cdot c_1(\sE)^2\cdot [H]^{n-2} \leq c_2(\sE)\cdot [H]^{n-2}.
\end{equation}
In case of equality, the sheaf $\sE$ is known to be projectively flat. We refer
the reader to Section~\vref{sec:3-1} for a brief discussion of projective
flatness.

Motivated by the structure theory of higher-dimensional projective manifolds, we
are particularly interested in the case where $\sE$ is the cotangent bundle $\Omega_X^1$
of an $n$-dimensional manifold $X$. In this setup, the equality case of
\eqref{eq:x1} reads
\begin{equation}\label{BGE}
\frac{n-1}{2n}\cdot c_1(X)^2\cdot [H]^{n-2} = c_2(X)\cdot [H]^{n-2}.
\end{equation}

While semistability of $\Omega^1_X$ occurs in many relevant cases and has important
geometric consequences, the Equality~\eqref{BGE} poses severe restrictions on
the geometry of~$X$.
\begin{itemize}
\item In case $K_X$ is ample, Equality~\eqref{BGE} will never hold, owing to the
stronger Miyaoka--Yau inequality.

\item If $K_X \equiv 0$, then by Yau's theorem, $X$ is an étale quotient of an
Abelian variety.

\item If $X$ is Fano and Kähler-Einstein, again the equality \eqref{BGE} cannot
occur, owing to the Chen--Oguie inequality.
\end{itemize}
In their remarkable paper \cite{MR3030068}, Jahnke and Radloff proved the
following complete characterisation of manifolds with semistable cotangent
bundle for which Equality~\eqref{BGE} holds.

\begin{thm}[\protect{Characterisation of torus quotients, \cite[Ths.\,0.1 \& 1.1]{MR3030068}}]
Let $X$ be a projective manifold of dimension $n$ and assume that $\Omega^1_X$ is
$H$-semistable for some ample line bundle $H$. If Equality~\eqref{BGE} holds,
then $X$ is a finite étale quotient of a torus. \qed
\end{thm}

In particular, this deals with manifolds of intermediate Kodaira dimension, and
moreover implies that in the Fano case Equality~\eqref{BGE} can never happen,
independent of the existence question for Kähler-Einstein metrics.

\subsection{Main result of this paper}

We have learned from numerous previous results, including \cite{GKP13, LT18,
GKT16, GKPT19, GKPT19b, GKP20a}, that the natural context for uniformisation
results is that of minimal model theory. The aim of this paper is therefore to
generalise the theorem of Jahnke--Radloff to the case when $X$ has klt
singularities.

\begin{thm}[Characterisation of quasi-Abelian varieties]\label{thm:charQAbelian}
Let $X$ be a projective klt space of dimension $n \geq 2$ and let $H \in \Div(X)$
be ample. Assume that $\Omega^{[1]}_X$ is semistable with respect to $H$ and that
its $\bQ$-Chern classes satisfy the equation
\begin{equation}\label{eq:qcce}
\frac{n-1}{2n}\cdot \what{c}_1 \bigl(\Omega^{[1]}_X \bigr)^2\cdot [H]^{n-2}
= \what{c}_2 \bigl(\Omega^{[1]}_X \bigr)\cdot [H]^{n-2}.
\end{equation}
Then, $X$ is quasi-Abelian and has at worst quotient singularities.
\end{thm}

In Theorem~\ref{thm:charQAbelian}, we say that a normal projective variety $X$
is \emph{quasi-Abelian} if there exists a quasi-étale cover $\wtilde{X} \to X$
from an Abelian variety $\wtilde{X}$ to $X$. The symbols $\what{c}_\sbullet$ denote
$\bQ$-Chern classes on the klt variety $X$, as recalled in
\cite[\S 3]{GKPT19b}.

\subsection{Normalised cotangent sheaves}

Work of Narasimhan, Seshadri and others, summarised for example in
\cite[Th.\,1.1]{MR3030068} and explained in detail by Nakayama in
\cite[Th.\,A]{NakayamaNormalizedPreprint}, can be used to reformulate
Jahnke-Radloff's result in terms of positivity properties of natural tensor
sheaves: a projective manifold $X$ of dimension $n$ is quasi-Abelian if and only
if the \emph{normalised cotangent bundle}, $\Sym^n \Omega^1_X \otimes \sO_X(-K_X)$, is nef.

While the arguments presented by Nakayama use intersection theory computations
on the total space of the projectivised cotangent bundle that cannot immediately
be carried over to singular varieties, we are nevertheless able to obtain the
analogous result in our setup.

\begin{thm}\label{thm:6-1}
Let $X$ be a normal projective variety of dimension $n \geq 2$. Assume that $X$
is klt and that the reflexive normalised cotangent sheaf
\[
\bigl(\Sym^n\Omega^1_X \otimes \sO_X(-K_X)\bigr)^{**}
\]
is nef. Then, $X$ is quasi-Abelian.
\end{thm}

Definition~\vref{def:possheaf} recalls the meaning of \emph{nef} for a sheaf
that is not necessarily locally free.

\subsection{Strategy of proof}

While the general strategy of proof is similar to the one employed by Jahnke and
Radloff, our argument introduces a number of new tools; this includes a detailed
analysis of sheaves that are projectively flat on the smooth a klt variety,
especially as their behaviour near the singularities is concerned. On the one
hand, these tools allow us to deal with the (serious) complications arising from
the singularities. On the other hand, they enable us to streamline parts of
Jahnke-Radloff's proof, thereby clarifying the underlying geometric principles.

As an intermediate step, we obtain the following result, which might be of
independent interest to some readers.

\begin{thm}[\protect{= Theorem~\ref{thm:druel} in Section~\ref{sec:4-2}}]
Let $X$ be a normal, projective klt variety of dimension $n \geq 2$. Assume that
the sheaf $\Omega^{[1]}_X$ of reflexive differentials is of the form
$\Omega^{[1]}_X \cong \sL^{\oplus n}$, where $\sL$ is reflexive of rank one. Then, $X$ is
quasi-abelian.
\end{thm}

\subsubsection*{Acknowledgements}

We thank Indranil Biswas and Stefan Schröer for answering our questions. We
thank the anonymous referee for an extremely helpful and detailed report that
improved the presentation of this paper considerably. Section~\ref{sec:4} has
been re-written, strengthened and substantially shortened in response to the
reviewer's comments. In particular, the final version of some of the results
and proofs in Section~\ref{sec:4-2} were suggested by the referee.

\section{Conventions, notation, and variations of standard facts}

\subsection{Global conventions}

Throughout this paper, all schemes, varieties and morphisms will be defined over
the complex number field. We follow the notation and conventions of
Hartshorne's book \cite{Ha77}. In particular, varieties are always assumed to
be irreducible. For all notation around Mori theory, such as klt spaces and klt
pairs, we refer the reader to \cite{KM98}.

\subsection{Varieties and complex spaces}

In order to keep notation simple, we will sometimes, when there is no danger of
confusion, not distinguish between algebraic varieties and their underlying
complex spaces. Along these lines, if $X$ is a quasi-projective complex
variety, we write $\pi_1(X)$ for the fundamental group of the associated complex
space.

\subsection{Reflexive sheaves}

As in most other papers on the subject, we will frequently consider reflexive
sheaves and take reflexive hulls. Given a normal, quasi-projective variety (or
normal, irreducible complex space) $X$, we write
$\Omega^{[p]}_X := \bigl(\Omega^p_X \bigr)^{**}$ and refer to this sheaf as \emph{the
sheaf of reflexive differentials}. More generally, given any coherent sheaf
$\sE$ on $X$, write $\sE^{[\otimes m]} := \bigl(\sE^{\otimes m} \bigr)^{**}$ and
$\det \sE := \bigl(\wedge^{\rank \sE} \sE \bigr)^{**}$. Given any morphism $f : Y \to X$ of
normal, quasi-projective varieties (or normal, irreducible, complex spaces), we
write $f^{[*]} \sE := (f^* \sE)^{**}$.

\subsection{Nef sheaves}

We recall the notion of a nef sheaf in brief and collect basic properties.

\begin{defn}[Nef and ample sheaves, \cite{MR670921}]\label{def:possheaf}
Let $X$ be a normal, projective variety and let $\sS \not = 0$ be a non-trivial
coherent sheaf on $X$. We call $\sS$ \emph{ample}/\emph{nef} if the locally
free sheaf $\sO_{\bP(\sS)}(1) \in \Pic(\bP(\sS))$ is ample/nef.
\end{defn}

We refer the reader to \cite{GrCa60} for the definition of $\bP(\sS)$, and to
\cite[\S 2 \&\,Th.\,2.9]{MR670921} for a more detailed discussion of
amplitude and for further references. We~mention a few elementary facts without
proof.

\begin{fact}[Nef sheaves]\label{fact:1}
Let $X$ be a normal, projective variety.
\begin{enumeraten}
\item\label{il:f1-1} A direct sum of sheaves on $X$ is nef iff every summand
is nef.
\item\label{il:f1-2} Pull-backs and quotients of nef sheaves are nef.
\item\label{il:f1-3} A sheaf $\sE$ is nef on $X$ if and only if for every smooth
curve $C$ and every morphism $\gamma: C \to X$, the pull-back $\gamma^* \sE$ is nef. \qed
\end{enumeraten}
\end{fact}

\subsection{Flat sheaves}

One key notion in our argument is that of a flat sheaf. We briefly recall the
definition.

\begin{defn}[\protect{Flat sheaf, \cite[Def.\,1.15]{GKP13}}]\label{defn:flat}
If $X$ is any normal, irreducible complex space and $\sF$ is any locally free
coherent sheaf on $X$, we call $\sF$ flat if it is defined by a
(finite-dimensional) complex representation of the fundamental group $\pi_1(X)$.
A locally free coherent sheaf on a quasi-projective variety is called flat if
the associated analytic sheaf on the underlying complex space is flat.
\end{defn}

\begin{rem}[Simple properties of flat sheaves]\label{rem:2-3}
Tensor powers, duals, symmetric products and wedge products of locally free,
flat sheaves are locally free and flat. The pull-back of a locally free, flat
sheaf under an arbitrary morphism is locally free and flat. If $\sF$ is a
locally free and flat sheaf on a normal, irreducible complex space $X$, then
there exists a description of $\sF$ in terms of a trivialising covering and
transition functions where all transitions functions are constant.
\end{rem}

The following lemma is a direct consequence of \cite[Chap.\,II,
Prop.\,3.1]{Kob87}. We leave the details to the reader.

\begin{lem}[Chern class of flat sheaf]\label{lem:c1offlat}
Every invertible, flat sheaf on a normal projective variety is numerically
trivial. \qed
\end{lem}

\subsection{Covering maps and quasi-étale morphisms}

A \emph{cover} or \emph{covering map} is a finite, surjective morphism
$\gamma : X \to Y$ of normal, quasi-projective varieties (or normal, irreducible
complex spaces). The covering map $\gamma$ is called \emph{Galois} if there exists a
finite group $G \subset \Aut(X)$ such that $\gamma$ is isomorphic to the quotient map.

A morphism $f : X \to Y$ between normal varieties (or normal, irreducible complex
spaces) is called \emph{quasi-étale} if $f$ is of relative dimension zero and
étale in codimension one. In other words, $f$ is quasi-étale if
$\dim X = \dim Y$ and if there exists a closed, subset $Z \subseteq X$ of codimension
$\codim_X Z \geq 2$ such that $f|_{X \setminus Z} : X \setminus Z \to Y$ is étale.

\subsection{Maximally quasi-étale spaces}
\label{ssec:mqes}

Let $X$ be a normal, quasi-projective variety (or a normal, irreducible complex
space). We say that $X$ is \emph{maximally quasi-étale} if the natural
push-forward map of fundamental groups,
\[
\pi_1(X_{\reg}) \To{(\text{incl})_*} \pi_1(X)
\]
induces an isomorphism between the profinite completions,
$\what{\pi}_1(X_{\reg}) \cong \what{\pi}_1(X)$.

\begin{rem}[Surjectivity of $(\text{incl})_*$]
Recall from \cite[0.7.B on p.\,33]{FL81} that the natural push-forward map
$(\text{incl})_*$ is always surjective.
\end{rem}

\begin{rem}[Existence of maximally quasi-étale covers]
If $(X,\Delta)$ is any quasi-projective klt pair, then $X$ admits a quasi-étale
cover $\gamma : \what{X} \to X$ where $\what{X}$ is maximally quasi-étale,
\cite[Th.\,1.14]{GKP13}. The pair $(\what{X}, \gamma^*\Delta)$ is again klt.
\end{rem}

\begin{rem}[Extension of representations and flat sheaves]\label{rem:extrep}
Assume that $X$ is a normal, irreducible complex space that is maximally
quasi-étale. Using Malcev's theorem, it has been shown in
\cite[\S 8.1]{GKP13} that any representation $\rho : \pi_1(X_{\reg}) \to \GL(N,\bC)$
factorises via the fundamental group of $X$,
\[
\begin{tikzcd}[column sep=2cm]
\pi_1(X_{\reg}) \ar[r, two heads, "\ts(\text{incl})_*"'] \ar[rr, bend left=15, "\ts\rho"] & \pi_1(X) \ar[r, "\ts\exists !"'] & \GL(N,\bC).
\end{tikzcd}
\]
In particular, it follows that any flat bundle on $X_{\reg}$ extends to a flat
bundle on $X$.
\end{rem}

\subsection{Local fundamental groups of contraction morphisms}

The behaviour of the fundamental group under extremal contractions and
resolutions of singularities has been studied by Takayama in a series of papers.
The following result is due to him.

\begin{prop}[Local fundamental groups of contraction morphisms]\label{prop:takayama}
Let $X$ be a projective klt variety and $f : X \to Y$ the contraction of a
$K_X$-negative extremal ray. If $y \in Y$ is any point, then there exists a
neighbourhood $U = U(y) \subseteq Y$, open in the analytic topology, such that
$f^{-1}(U)$ is connected and simply connected.
\end{prop}
\begin{proof}
Takayama formulates a global version of the result in
\cite[Th.\,1.2]{Takayama2003}. The arguments of \cite[Proof of Th.\,1.2 on
p.\,834]{Takayama2003} apply in our setup and give a full proof of
Proposition~\ref{prop:takayama}.
\end{proof}

\subsection{Abundance}

We will later use the following special case of the abundance conjecture. Even
though the proof uses only standard results from the literature, we found it
worth to include it here in full, for the reader's convenience and for later
referencing.

\begin{prop}\label{prop:abundance}
Let $f : X \dasharrow Y$ be a dominant, almost holomorphic, rational map with
connected fibres between projective varieties. Assume that
\begin{enumeraten}
\item\label{il:AS1} The variety $Y$ is smooth, positive-dimensional and of
general type.

\item\label{il:AS2} The variety $X$ is klt, and the canonical divisor $K_X$ is
nef.

\item\label{il:AS3} The variety $X$ is smooth around general fibres of $f$.

\item\label{il:AS4} The general fibre is smooth and is a good minimal model.
\end{enumeraten}
Then, $X$ is a good minimal model. In other words, $K_X$ is semiample.
\end{prop}
\begin{proof}
Blowing up $X$ in the indeterminacy locus, we obtain a diagram
\[
\begin{tikzcd}
\wtilde{X} \ar[r, "\ts\pi"] \ar[d, "\ts\wtilde{f}"'] & X \ar[d, dashed, "\ts f"] \\
Y \ar[r, equal] & Y
\end{tikzcd}
\]
where $\pi$ is isomorphic around the general fibre $X_y \cong \wtilde{X}_y$ of $f$
or $\wtilde{f}$, respectively. Next, consider the relative Iitaka fibration
for $\pi^* K_X$ on $\wtilde{X}/Y$. As discussed in \cite[\S 2.4]{MR3286535},
this is given by the (standard) Iitaka fibration for a line bundle of the form
\[
\sL := \pi^* K_X + \wtilde{f}^*(\text{sufficiently ample}).
\]
Since sufficiently high powers of $\sL$ have no base points along the general
fibre $\wtilde{X}_y$, we may replace $\wtilde{X}$ by a further blow-up and
obtain a diagram as follows,
\[
\begin{tikzcd}
\wtilde{X} \ar[r, equal] \ar[d, "\ts\text{rel.~Iitaka}"'] & \wtilde{X} \ar[r, "\ts\pi"] \ar[d, "\ts\wtilde{f}"'] & X \ar[d, dashed, "\ts f"] \\
Z \ar[r] & Y \ar[r, equal] & Y
\end{tikzcd}
\]
where $Z$ is smooth and where $\sL$ is numerically trivial on the general fibre
$\wtilde{X}_z$ of $\wtilde{X}$ over $Z$. The numerical dimension $\nu(X)$ and
the Kodaira dimension $\kappa(X)$ are then estimated in terms of $\wtilde{X}_y$ and
$\dim Y$ as follows,
\begin{align*}
\kappa(\wtilde{X}_y) + \dim Y & = \dim Z && \text{\cite[\S 2.4]{MR3286535}} \\
& = \nu(\pi^*K_X|_{\wtilde{X}_z}) + \dim Z && \text{since } \pi^*K_X|_{\wtilde{X}_z} \equiv 0\\
& \geq \nu(\pi^* K_X) = \nu(K_X) && \text{\cite[V.~Lem.\,2.3.(2)]{Nakayama04}} \\
& \geq \kappa(K_X) && \text{\cite[Prop.\,2.2]{MR782236}} \\
& \geq \kappa(\wtilde{X}) && \text{because $\pi$ is birational}\\
& \geq \kappa(\wtilde{X}_y) + \dim Y && \text{by \cite[Satz~III]{MR641815}}.
\end{align*}
It follows that $\kappa(X) = \nu(X)$, and hence by \cite[Cor.\,6-1-13]{KMM87} that
$K_X$ is semiample, as desired.
\end{proof}

\subsection{The negativity lemma}

For later reference, we note the following two minor generalisations of the
classic negativity lemma, as formulated for instance in
\cite[Lem.\,3.39--3.41]{KM98}.

\begin{lem}[Negativity Lemma I]\label{lem:negativity1}
Let $f : X \to Y$ be a surjective morphism of normal, projective varieties. If
$E \in \bQ\Div(X)$ is non-zero and effective with $\codim_Y f(\supp E) \geq 2$, then
$E$ is not nef.
\end{lem}
\begin{proof}
Choose a very ample line bundle $\sL$ on $X$ and a general tuple of sections
$(H_1, \dots, H_{\dim X-\dim Y}) \in |\sL|^{\times (\dim X-\dim Y)}$. Set
$S := H_1 \cap\cdots\cap H_{\dim X-\dim Y}$. This is a normal subvariety of $X$. The
restricted morphism $f|_S$ is surjective and generically finite, and $E \cap S$
is a non-zero, effective, $\bQ$-Cartier $\bQ$-divisor on $S$. The Stein
factorisation of $f|_S$ is then birational and contracts $E \cap S$. The
negativity lemma, \cite[Lem.\,3.39(1)]{KM98} with $-B = E \cap S$, then applies to
conclude that $E \cap S$ (and therefore~$E$) are non-nef.
\end{proof}

\begin{lem}[Negativity Lemma II]\label{lem:negativity2}
Let $f : X \to Y$ be a surjective morphism of normal, projective varieties, where
$Y$ is a curve. If $E \in \bQ\Div(X)$ is non-zero, effective, and maps to a point
in $Y$, then either $\kappa(E) = 1$ or $E$ is not nef.
\end{lem}
\begin{proof}
Write $y := f(E)$. If $\supp(E)$ equals the set-theoretic fibre $f^{-1}(y)$,
then $\kappa(E) = 1$. We will therefore assume that $\supp(E)$ is a proper subset
of $f^{-1}(y)$. Choose a very ample line bundle $\sL$ on $X$ and a general
tuple of sections $(H_1, \dots, H_{\dim X-2}) \in |\sL|^{\times (\dim X-\dim Y)}$. Set
$S := H_1 \cap\cdots\cap H_{\dim X-2}$. This is a normal surface in $X$. The divisor
$E \cap S$ is a non-zero, effective, $\bQ$-Cartier $\bQ$-divisor on $S$, supported on
a proper subset of $(f|_S)^{-1}(y)$. Zariski's Lemma, \cite[Chap.\,III,
Lem.\,8.2]{HBPV}, will then show that $(E \cap S)^2 < 0$. It follows that $E \cap S$
(and therefore $E$) cannot be nef.
\end{proof}

\subsection{Abelian group schemes}

The proof of our main result, the characterisation of quasi-Abelian varieties in
Theorem~\ref{thm:charQAbelian}, uses a minor generalisation of Kollár's
characterisation of étale quotients of Abelian group schemes,
\cite[Th.\,6.3]{Kollar93}. We recall the relevant notions first.

\begin{defn}[Abelian group scheme]
An \emph{Abelian group scheme over a base $B$} is a smooth, proper morphism
$a: A \to B$ between smooth varieties such that every fibre of $a$ is an Abelian
variety, and such that there exists a section $B \to A$.
\end{defn}

\begin{defn}[\protect{Generically large fundamental group, \cite[Defn.~6.1]{Kollar93}}]
Let $X$ be a normal, projective variety and let $Y \subset X$ be a closed
subvariety. We say that \emph{$X$ has a generically large fundamental group
on $Y$} if for all very general points $y \in Y$ and for every closed and
positive-dimensional subvariety $y \in Z \subset Y$ with normalisation $\overline{Z}$,
the image of the natural morphism $\pi_1\bigl(\overline{Z} \bigr) \to \pi_1(X)$ is
infinite.
\end{defn}

The generalisation of Kollár's result is then formulated as follows.

\begin{prop}[Characterisation of étale quotients of Abelian group schemes]\label{prop:K63}
Let $f : X \to Y$ be a surjective morphism with connected fibres between normal,
projective varieties. Assume that there exists $\Delta \in \bQ\Div(X)$ such that
$(X,\Delta)$ is klt. Let $y \in Y$ be a very general point with fibre $X_y$ and
assume that the following holds.
\begin{itemize}
\item The fibre $X_y$ has a finite étale cover that is birational to an
Abelian variety.

\item The variety $X$ has a generically large fundamental group on $X_y$.
\end{itemize}
Then, there exists an étale cover $\gamma_X : \what{X} \to X$ such that the fibration
$\what{a} : \what{A} \to \what{Y}$ obtained as the Stein factorisation of
$(f\circ \gamma_X) : \what{X} \to Y$ is birational to an Abelian group scheme over a
proper base.
\end{prop}
\begin{proof}
If $X$ and $Y$ are smooth, this is \cite[Th.\,6.3]{MR1341589}. If not, choose
$\Delta \in \bQ\Div(X)$ such that $(X,\Delta)$ is klt. We will consider resolutions of the
singularities and construct a diagram of surjective morphisms between normal,
projective varieties,
\begin{equation}\label{eq:D2a}
\begin{tikzcd}[column sep=2.5cm]
\wcheck{X} \ar[r, "\gamma_{\wtilde{X}}\text{, finite and étale}"] \ar[d, "\wcheck{a} \text{, conn.~fibres}"'] & \wtilde{X} \ar[r, "\pi_X \text{, resolution}"] \ar[d, "\wtilde{f} \text{, conn.~fibres}"] & X \ar[d, "f \text{, conn.~fibres}"] \\
\wcheck{Y} \ar[r, "\gamma_{\wtilde{Y}}\text{, finite}"'] & \wtilde{Y} \ar[r, "\pi_Y \text{, resolution}"'] & Y
\end{tikzcd}
\end{equation}
where the morphism $\wcheck{a}$ is birational to an Abelian group scheme over
a proper base. For this, begin by choosing appropriate resolution morphisms
$\pi_X$ and $\pi_Y$ to construct the right square in \eqref{eq:D2a}. Once this is
done, choose a very general point $\wtilde{y} \in \wtilde{Y}$ with very general
image point $y := \pi_Y(\wtilde{y})$. Using a result of Campana,
\cite[Prop.\,1.3]{Campana91}, we see that the variety $\wtilde{X}$ has a
generically large fundamental group on $\wtilde{X}_{\wtilde{y}}$. Kollár's
result therefore applies to $\wtilde{f}$ and gives the left square.

There is more: using the assumption that $(X,\Delta)$ is klt, Takayama has shown in
\cite[Th.\,1.1]{Takayama2003} that the natural push-forward map
$(\pi_X)_*: \pi_1\bigl(\wtilde{X}\,\bigr) \to \pi_1(X)$ is isomorphic. The étale
morphism $\gamma_{\wtilde{X}}$ therefore comes from an étale cover of $X$, say
$\gamma_X : \what{X} \to X$. Stein factorising $\pi_X\circ \gamma_{\wtilde{X}}$ and
$\pi_Y\circ \gamma_{\wtilde{Y}}$, we will therefore obtain a diagram of the following
form,
\begin{equation}\label{eq:D2b}
\begin{tikzcd}[column sep=2.5cm]
\wcheck{X} \ar[rr, "\pi_X\circ \gamma_{\wtilde{X}}", bend left=15] \ar[r, "\pi_{\what{X}}\text{, resolution}"'] \ar[d, "\wcheck{a} \text{, conn.~fibres}"'] & \what{X} \ar[r, "\gamma_X \text{, étale}"'] \ar[d, "\what{a}", dashed] & X \ar[d, "f \text{, conn.~fibres}"] \\
\wcheck{Y} \ar[rr, "\pi_Y\circ \gamma_{\wtilde{Y}}"', bend right=15] \ar[r, "\pi_{\wtilde{Y}}\text{, conn.~fibres}"] & \what{Y} \ar[r, "\gamma_Y \text{, finite}"] & Y
\end{tikzcd}
\end{equation}
where $\what{a} := \pi_{\wtilde{Y}}\circ \wcheck{a}\circ \pi^{-1}_{\what{X}}$. This map
makes \eqref{eq:D2b} commute and is birational to $\wcheck{a}$, which is in
turn birational to an Abelian group scheme. To end the proof, it will
therefore suffice to show that $\what{a}$ is actually a morphism. This is not
so hard: if $\what{x} \in \what{X}$ is any point with preimage
$Z := \pi^{-1}_{\what{X}} \bigl(\what{x}\, \bigr)$, then
\[
\bigl(\gamma_Y\circ \pi_{\wtilde{Y}} \circ \wcheck{a}\,\bigr)(Z) =
f\bigl(\gamma_X(\what{x}\,)\bigr) = \text{point in } Y
\]
but since $\gamma_Y$ is finite, this implies that
$\bigl(\pi_{\wtilde{Y}} \circ \wcheck{a}\,\bigr)(Z)$ is a point in $\what{Y}$. The
Rigidity Lemma in the form of \cite[Lem.\,1.15]{Debarre01} will then show that
$\what{a}$ is well-defined at~$\what{x}$.
\end{proof}

\section{Projective flatness}

\subsection{Projectively flat bundles and sheaves}
\label{sec:3-1}

\emph{Projective flatness} is the core notion of this paper. We will only
recall the most relevant definition here and refer the reader to
\cite[\S 3]{GKP20a} for a detailed discussion of projectively flat bundles
and sheaves.

\begin{defn}[Projectively flat bundles and sheaves on complex spaces]\label{def:3-2}
Let $X$ be a normal and irreducible complex space, let $r \in \bN$ be any number
and let $\bP \to X$ be a $\bP^r$-bundle. We call the bundle $\bP \to X$
\emph{(holomorphically) projectively flat} if there exists a representation of
the fundamental group, $\rho : \pi_1(X) \to \bP\negthinspace\GL(r+1,\bC)$, and an
isomorphism of complex spaces over $X$
\[
\bP \cong_X \factor{\wtilde{X} \times \bP^r}{\pi_1(X)},
\]
where $\wtilde{X}$ is the universal cover of $X$ and where the action
$\pi_1(X) \circlearrowleft \wtilde{X} \times \bP^r$ is the diagonal action. A locally free sheaf $\sF$
of $\sO_X$-modules is called \emph{(holomorphically) projectively flat} if the
associated bundle $\bP(\sF)$ is projectively flat.
\end{defn}

\begin{defn}[Projectively flat bundles and sheaves on complex varieties]
Let $X$ be a connected, complex, quasi-projective variety and let $r \in \bN$ be
any number. A $\bP^r$-bundle $\bP \to X$ is called \emph{projectively flat} if the
associated analytic bundle $\bP^{(\an)} \to X^{(\an)}$ is projectively flat. Ditto
for coherent sheaves.
\end{defn}

\subsection{Projectively flat bundles induced by a solvable representation}

One way to analyse a projectively flat sheaf $\sE$ on a manifold or variety $X$ is
to look at the Albanese map $X \to \Alb(X)$, and to try and understand the
restriction $\sE|_F$ to a general fibre. The key insight here is that the
Albanese map is a Shafarevich map for the commutator subgroup,
\cite[Prop.\,4.3]{Kollar93}. This allows to relate the geometry of the Albanese
map to the group theory of the representation that defines the projectively flat
structure on $\sE$. The following proposition, which is a variation of material
that appears in \cite{MR3030068}, is key to our analysis.

\begin{prop}[Projectively flat bundles induced by a solvable representation]\label{prop:pfbisr}
Let $X$ be a normal, irreducible complex space that is maximally quasi-étale.
Let $\sF$ be a rank-{}$n$ reflexive sheaf on $X$ such that $\sF|_{X_{\reg}}$ is
locally free and projectively flat, so that $\bP(\sF|_{X_{\reg}})$ is isomorphic
to a projective bundle $\bP_\rho$, for a representation
$\rho : \pi_1(X_{\reg}) \to \PGL(n,\bC)$. If the image of $\rho$ is contained in a
connected, solvable, algebraic subgroup $G \subset \PGL(n,\bC)$, then there exists an
isomorphism
\begin{equation}\label{eq:OmegaFA0-1}
\sF \cong \sF' \otimes \sA,
\end{equation}
where $\sA$ is Weil divisorial\footnote{Weil divisorial sheaf = reflexive sheaf of rank one.} and where $\sF'$ is locally free and admits a filtration
\begin{equation}\label{eq:3-3-2}
0 = \sF'_0 \subsetneq \sF'_1 \subsetneq \cdots \subsetneq \sF'_n = \sF'
\end{equation}
such that the following holds.
\begin{enumeraten}
\item\label{il:x1} All sheaves $\sF'_i$ are locally free and flat in the sense
of Definition~\ref{defn:flat}, with $\rank (\sF'_i)$ equal to $i$.

\item\label{il:x2} The sheaf $\sF'_1$ is trivial, $\sF'_1 \cong \sO_X$.
\end{enumeraten}
\end{prop}

\begin{rem}\label{rem:pfbisr}
In the setting of Proposition~\ref{prop:pfbisr}, note that there exists an
isomorphism of Weil divisorial sheaves, $\det \sF \cong \sA^{[\otimes n]} \otimes \det \sF'$.
Recall from Remark~\ref{rem:2-3} that the invertible factor $\det \sF'$ is flat
and hence by Lemma~\ref{lem:c1offlat} numerically trivial.
\end{rem}

\begin{proof}[Proof of Proposition~\ref{prop:pfbisr}]
To begin, we claim that the projective representation $\rho$ lifts to a linear
representation, which we will denote by $\sigma$. To formulate our result more
precisely, write $B(n,\bC) \subset \GL(n,\bC)$ for the subgroup of upper-triangular
matrices and $B_1(n,\bC) \subset B(n,\bC)$ for the subgroup of upper-triangular matrices
whose top-left entry equals one. In other words,
\[\arraycolsep4.5pt
B_1(n,\bC) := \left\{ \left. \:\:
\begin{pNiceMatrix}[nullify-dots]
1 & * & \Cdots & & * \\
0 & * & \Ddots & & \Vdots \\
\Vdots & \Ddots & \Ddots & & \\
& & & & * \\
0 & \Cdots & & 0 & *
\end{pNiceMatrix}
\:\:\right|\:\: \text{all } * \in \bC
\right\} = \Bigl\{ (b_{\sbullet\sbullet}) \in B(n, \bC) \mid b_{11}=1\Bigr\}.
\]
With this notation, we claim there exists a connected linear group
$\what{G} \subseteq B_1(n,\bC) \subset \GL(n,\bC)$ such that the natural projection
$\pi : \GL(n,\bC) \to \PGL(n,\bC)$ induces an isomorphism between $\what{G}$ and $G$.
To this end, let $\what{G}' \subset \SL(n, \bC)$ be the maximal connected subgroup of
$\pi^{-1}(G) \cap \SL(n, \bC)$. The restricted morphism
$\pi|_{\what{G}'} : \what{G}' \to G$ is then finite and surjective, and induces an
isomorphism of the Lie algebras $\Lie(\what{G}')$ and $\Lie(G)$. In
particular, $\what{G}'$ is solvable. Therefore, choosing an appropriate basis
of $\bC^{n}$, Lie's Theorem allows us to assume without loss of generality that
$\what{G}'$ is contained in the subgroup $B(n,\bC)$ of upper-triangular
matrices. Finally, set $\what{G} := \pi^{-1}(G) \cap B_1(n, \bC)$. The obvious fact
that any invertible upper triangular matrix $A \in \what{G}' \subseteq B(n,\bC)$ has a
unique multiple $\lambda\cdot A$ that lies in $B_1(n, \bC)$ implies that the restriction
$\pi|_{\what{G}} : \what{G} \to G$ is then isomorphic, as desired.

Remark~\ref{rem:extrep} and the assumption that $X$ is maximally quasi-étale
allow us to identify representations of $\pi_1(X)$ and $\pi_1(X_{\reg})$. With
these identifications made, the following commutative diagram summarises the
situation,
\[
\begin{tikzcd}[column sep=large]
& \what{G} \arrow[r, hook] \ar[d, "\ts\cong"] & B_1(n,\bC) \arrow[r, hook] & \GL(n,\bC) \arrow[d, "\ts\pi\text{, nat.~projection}"] \\
\pi_1(X) \arrow[r, "\ts\rho"'] \arrow[ur, "\ts\sigma", bend left=20] & G \arrow[rr, hook] & & \PGL(n,\bC)
\end{tikzcd}
\]

If $\sF'$ denotes the locally free, flat sheaf on $X$ associated with $\sigma$, then
the factorisation of $\sigma$ via the group $B_1(n, \bC)$ implies that $\sF'$ admits a
filtration as in \eqref{eq:3-3-2} that satisfies conditions \ref{il:x1} and
\ref{il:x2}. By construction, we have $\bP(\sF|_{X_{\reg}}) \cong \bP(\sF'|_{X_{\reg}})$,
and so there exists an invertible sheaf $\sA_{\reg}$ on $X_{\reg}$ such that
$\sF|_{X_{\reg}} \cong \sF'|_{X_{\reg}} \otimes \sA_{\reg}$. A presentation of $\sF$ as in
Equation~\eqref{eq:OmegaFA0-1} therefore exists by taking $\sA := \iota_* \sA_{\reg}$,
where $\iota : X_{\reg} \to X$ is the inclusion.
\end{proof}

\section{Varieties with splitting cotangent sheaves}
\label{sec:4}

One relevant case that keeps appearing in the discussion of varieties with
projectively flat sheaf of reflexive differentials is the one where $\Omega^{[1]}_X$
decomposes as a direct sum of Weil divisorial sheaves. We discuss settings
where there are local or global decompositions of this form. The results of
this section will be used later in the proof of our main result, but might be of
independent interest. In relation to \cite{MR3030068}, the local results are
new, prompted by the appearance of singularities in our context, and the results
on varieties with globally split cotangent sheaf allow us to replace some of the
more involved arguments of \loccit\ by a clearer structure result.

\subsection{Varieties where $\Omega^{[1]}_X$ is locally split}
\label{sect:split}

Given a variety $X$ where $\Omega^{[1]}_X$ is projectively flat, for instance because
the flatness criterion of \cite[Th.\,1.6]{GKP20a} applies, the local description
of projectively flat sheaves, \cite[Prop.\,3.11]{GKP20a}, can often be used to
reduce to the case where $\Omega^{[1]}_X$ decomposes locally. We will show that this
already forces $X$ to have quotient singularities.

\begin{prop}\label{prop:icq}
Let $X$ be a normal complex space with klt singularities. Assume that the
sheaf $\Omega^{[1]}_X$ of reflexive differentials is of the form
\begin{equation}\label{eq:icq}
\Omega^{[1]}_X \cong \sL^{\oplus \dim X},
\end{equation}
where $\sL$ is reflexive of rank one. Then, $X$ has only isolated, cyclic
quotient singularities.
\end{prop}

We prove Proposition~\ref{prop:icq} in Section~\vref{ssec:popicq}.

\begin{rem}
If $X$ is a quasi-projective variety with klt singularities where $\Omega^{[1]}_X$
is of the form $\sL^{\oplus \dim X}$ for a Weil-divisorial sheaf $\sL$, then the proof
of Proposition~\ref{prop:icq} applies verbatim to show that $X$ is
Zariski-locally a quotient of a smooth variety. We do not use this fact
however.
\end{rem}

\subsection{Varieties where $\Omega^{[1]}_X$ is globally split}
\label{sec:4-2}

There are settings where $\Omega^{[1]}_X$ decomposes globally, for instance because
it is projectively flat and $X$ is known to be simply connected. This obviously
puts strong conditions on the global geometry of $X$; see \cite{MR1758581} for
further results in the smooth setup.

The following proposition, including the proof, is due to the referee. He also
suggested Theorem~\ref{thm:druel}, which is stronger than our original result
that included a pseudo-effectivity assumption. Our previous arguments were of
analytic nature and used a theorem of Demailly~\cite{Dem02}.

\begin{prop}\label{prop:druel}
Let $U$ be a connected complex manifold of dimension $n \geq 2$. Assume that
$\Omega^1_U \cong \sL^{\oplus n}$, where $\sL$ is invertible. Then,
$c_1(\sL) = 0 \in H^1 \bigl(U,\, \Omega^1_U \bigr)$.
\end{prop}
\begin{proof}
Consider the induced decomposition on cohomology
\[
H^1 \bigl(U,\, \Omega^1_U\bigr) \cong H^1\bigl(U,\, \sL \bigr)^{\oplus n}.
\]
For reasons that
will become apparent in a second, we are particularly interested in the
coordinate hyperplanes
\[
Z_i := \left\{ (\alpha_1, \dots, \alpha_n) \in H^1\bigl(U,\, \sL \bigr)^{\oplus n} \bigm| \alpha_i = 0
\right\} \subsetneq H^1 \bigl(U,\, \Omega^1_U \bigr).
\]
To see how the $Z_i$ come into play, restate the assumption as a decomposition
of the tangent bundle: $\sT_U \cong \sA_1 \oplus \cdots \oplus \sA_n$. The summands $\sA_\sbullet \cong \sL^*$ are
regular rank-one foliations on $U$, with normal bundles
\[
\sN_i = \sA_1 \oplus \cdots \oplus \what{\sA_i} \oplus \cdots \oplus \sA_n.
\]
In this setting, a theorem of Baum-Bott, \cite[Prop.\,3.3 \&
Cor.~3.4]{MR0261635} but see also \cite[Lem.\,3.1]{MR1760872}, describes the
first Chern class $c_1(\sN_i)$ as lying in the subspace
\[
Z_i = H^1 \bigl(U,\, \sN_i^* \bigr) \subsetneq H^1 \bigl(U,\, \Omega^1_U \bigr).
\]
But since $c_1(\sL)$ is proportional to any of the $c_1(\sN_\sbullet)$, we find that
$c_1(\sL)$ is contained in $Z_1 \cap\cdots\cap Z_n = \{0\}$. The assertion is thus shown.
\end{proof}

Applying Proposition~\ref{prop:druel} to the smooth locus of a klt variety, we
obtain the following characterisation of quasi-abelian varieties.

\begin{thm}\label{thm:druel}
Let $X$ be a normal, projective klt variety of dimension $n \geq 2$. Assume that
the sheaf $\Omega^{[1]}_X$ of reflexive differentials is of the form
$\Omega^{[1]}_X \cong \sL^{\oplus n}$, where $\sL$ is reflexive of rank one. Then, $X$ is
quasi-abelian.
\end{thm}
\begin{proof}
We have seen in Proposition~\ref{prop:druel} that the $\bQ$-Cartier sheaves $\sL$
are numerically trivial, and then so is $\omega_X \cong \sL^{[\otimes n]}$. Since $X$ is klt,
it follows from \cite[Prop.\,V.4.9]{Nakayama04} or \cite[Th.\,4.2]{MR2134273}
that $\omega_X$ and $\sL$ are torsion. As a consequence, we find a finite
quasi-étale cover $\gamma: \what{X} \to X$ such that the reflexive pull-back
$\gamma^{[*]} \sL$ is trivial. The space $\what{X}$ is klt and its tangent sheaf
$\sT_{\what{X}} \cong \left(\gamma^{[*]} \sL \right)^{\oplus n}$ is likewise trivial. By the
solution of the Lipman-Zariski conjecture for spaces with klt singularities,
\cite[Th.\,6]{GKKP11}, it follows that the covering space $\what{X}$ is
smooth. In other words, $\what{X}$ is a parallelisable projective manifold.
By a classic result of Wang, $\what{X}$ is an Abelian variety; see for example
\cite[\S 3.9, Lem.\,1]{MR1334091}.
\end{proof}

\subsection{Proof of Proposition~\ref*{prop:icq}}
\label{ssec:popicq}

The proof of Proposition~\ref{prop:icq} uses the following lemma that might be
of independent interest.

\begin{lem}\label{lem:glcs}
Let $X$ be a normal, complex space with only cyclic quotient singularities.
Let $S \subseteq X_{\sing}$ be any irreducible component of the singular locus. Then,
there exists a non-empty, open set $S^\circ \subseteq S$ such that every $x \in S^\circ$ admits an
open neighbourhood $U=U(x) \subseteq X$ of the form $U = W \times V$, where $W$ is smooth
and $V$ has an isolated quotient singularity.
\end{lem}
\begin{proof}
To begin, consider the special case where $V$ is a complex vector space,
$\rho : \bZ_m \to \GL(V)$ is a linear representation of a cyclic group and where
$X \subseteq V / \bZ_m$ is an open neighbourhood of $[\vec 0]$. Recall that $\rho$
decomposes into direct a sum of one-dimensional representations. Identifying
$\bZ_m$ with roots of unity, we can thus find linear coordinates on $V$ such
that multiplication with $\xi \in \bZ_m$ is given as
\[
\xi\cdot (v_1, \dots, v_m) = \bigl(\xi^{n_1}\cdot v_1, \dots, \xi^{n_m}\cdot v_m\bigr).
\]
An elementary computation in these coordinates shows that $S^\circ$ exists, and can
in fact be chosen to be dense in $S$.

Returning to the general case where $X$ is arbitrary, let $s \in S$ be any
point. By~assumption, there exists an open neighbourhood $U = U(x)$ and a
cyclic cover $\gamma : \wtilde U \to U$, say with group $G$. Shrinking $U$ and
passing to a subgroup, if need be, we may assume without loss of generality
that $\gamma$ is totally branched over $s$, that is $\gamma^{-1}(s) = \{ \wtilde s \}$.
The cyclic group thus acts on $\wtilde U$ and fixes $\wtilde s$. The group
$G$ clearly acts on the vector space $T_{\wtilde U}|_{\wtilde s}$ and
linearisation at the fixed point, \cite[Proof of Th.\,4]{MR0084174} or
\cite[Cor.\,2 on p.\,13]{MR782881}, shows the existence of a $G$-invariant open
neighbourhood $\wtilde U'$ of $\vec 0 \in T_{\wtilde U}|_{\wtilde s}$ and an
equivariant, biholomorphic map between $\wtilde U$ and $\wtilde U$. In other
words, we are in the situation discussed in the first paragraph of this proof,
where the claim has already been shown.
\end{proof}

\begin{proof}[Proof of Proposition~\ref{prop:icq}]
As a first step in the proof of Proposition~\ref{prop:icq}, we claim that $\sL$
is $\bQ$-Cartier. In fact, taking determinants of both sides in \eqref{eq:icq},
we obtain that $\sL^{[\otimes \dim X]} \cong \omega_X$, which is $\bQ$-Cartier by assumption.
Choose a minimal number $N \in \bN^+$ such that $\sL^{[\otimes N]}$ is locally free. The
reflexive tensor product $\bigl(\Omega^{[1]}_U \bigr)^{[\otimes N]}$ is then likewise
locally free.

As a second step, we show that $X$ has cyclic quotient singularities. In
fact, given any $x \in X$, we find an open neighbourhood $U=U(x)$ over which
$\sL^{[\otimes N]}$ is trivial. Chose a trivialisation and let $\gamma: \wtilde{U} \to U$ be
the associated index-one cover, which is cyclic of order $N$. The complex
space $\wtilde{U}$ has again klt singularities and
$\gamma^{[*]}\, \sL \cong\nobreak \sO_{\wtilde{U}}$. In particular, it follows that
$\Omega^{[1]}_{\wtilde{U}} \cong \gamma^{[*]} \Omega^{[1]}_U$ and $\sT_{\wtilde U}$ are both free.
The solution of the Lipman-Zariski conjecture for spaces with klt
singularities, \cite[Th.\,3.8 \& comments after Th.\,1.1]{Druel13a} or
\cite[Th.\,1.13]{KS18}, then asserts that $\wtilde{U}$ is smooth. We conclude
that $X$ has cyclic quotient singularities only.

As a third and last step, we show that the singularities of $X$ are isolated.
Assume to the contrary and let $S \subseteq X_{\sing}$ be a positive-dimensional,
irreducible component. We have seen in Lemma~\ref{lem:glcs} that there a
non-empty, open subset $S^\circ \subseteq S$ such that every $x \in S^\circ$ admits an open
neighbourhood $U = U(x) \subseteq X$ of product form $U = W \times V$, where $W$ is smooth
and $V$ has an isolated quotient singularity. In particular,
$\dim W = \dim S$ and
\[
\Omega^{[1]}_U \cong \Omega^1_W \oplus \Omega^{[1]}_V.
\]
In particular, the reflexive tensor power $\bigl(\Omega^{[1]}_U \bigr)^{[\otimes N]}$ is
written as a direct sum of reflexive sheaves with
$\bigl(\Omega^{[1]}_W \bigr)^{[\otimes (N-1)]} \otimes \Omega^{[1]}_V$ as one of its direct
summands. It follows that $\bigl(\Omega^{[1]}_U \bigr)^{[\otimes N]}$ is not locally
free and accordingly, neither is $\bigl(\Omega^{[1]}_X \bigr)^{[\otimes N]}$. This
contradicts the results obtained in the first paragraph of this proof, and
therefore finishes the proof of Proposition~\ref{prop:icq}.
\end{proof}

\section{Proof of Theorem~\ref*{thm:charQAbelian}}

\subsection{Preparation for the proof}

Using the criterion for projective flatness spelled out in
\cite[Th.\,1.6]{GKP20a}, the assumptions made in Theorem~\ref{thm:charQAbelian}
imply that $\Omega^1_{X_{\reg}}$ is projectively flat, at least after going to a
quasi-étale cover. This allows to apply following lemma in various settings
that appear throughout the proof of Theorem~\ref{thm:charQAbelian}.

\begin{lem}[Consequences of projectively flat cotangent bundle]\label{lem:6-5}
Let $X$ be a smooth, quasi-projective variety of dimension $n$ where $\Omega^1_X$ is
projectively flat.
\begin{enumeraten}
\item\label{il:6-5-1} If $K_X$ is nef, then $X$ does not contain complete
rational curves.

\item\label{il:6-5-2} Assume that there exists an immersion $\eta : F \to X$ where
$F$ is projective and smooth. If $\eta^* K_X$ is nef and if $K_F \equiv 0$,
then $F$ is quasi-Abelian.
\end{enumeraten}
\end{lem}
\begin{proof}
We prove the statements separately. Both proofs use the fact that every
locally free, projectively flat sheaf on a simply connected space is
isomorphic to a direct sum of the form $\sL^{\oplus\,\sbullet}$ where $\sL$ is invertible,
cf.~\cite[Prop.\,3.11]{GKP20a}.

As to the first statement, let $\eta: \bP_1 \to X$ be a non-constant map. Since
$\eta^*\Omega^1_X$ is projectively flat, there is an integer $m$ such that
$\eta^*\Omega^1_X \cong \sO_{\bP_1}(m)^{\oplus n}$. The canonical morphism $\eta^*\Omega^1_X \to \Omega^1_{\bP_1}$
then yields $m \leq -2$, contradicting the nefness of $\eta^*K_X$.

As to the second statement, applying the classic Decomposition Theorem
\cite[Th.\,1]{Bea83} and possibly passing to a finite étale cover, we may
assume that $F$ is of product form, $F \cong A \times Z$, where $A$ is Abelian
(possibly a point) and where $Z$ is simply connected with $K_Z = 0$. We aim
to prove that $Z$ is a point, argue by contradiction and assume that
$m := \dim Z$ is positive. Let $\iota: Z \to F$ denote the inclusion morphism. The
pull-back $(\eta\circ \iota)^* \Omega^1_X$ is projectively flat and hence of the form $\sL^{\oplus n}$,
for a suitable line bundle $\sL \in \Pic(Z)$. The obvious surjection
\[
\begin{tikzcd}[column sep=1.2cm]
\sL^{\oplus n} \cong (\eta\circ \iota)^* \Omega^1_X \ar[r, two heads, "\ts\diff (\eta\circ \iota)"] & \Omega^1_Z
\end{tikzcd}
\]
induces a non-trivial map $\sL^{\otimes m} \to \omega_Z \cong \sO_Z$. But then the assumption that
$\eta^*K_X$ and hence also $\sL$ are nef shows that $\sL^{\otimes m}$ cannot be a proper
ideal sheaf, so $\sL^{\otimes m} \cong \sO_Z$. It follows that either $\sL$ is torsion or that
$H^0\bigl(Z,\, \Omega^1_Z\bigr) \ne 0$, both contradicting the simple connectedness
of $Z$.
\end{proof}

\subsection{Proof of Theorem~\ref*{thm:charQAbelian}}
\label{ssec:pf44}

We maintain notation and assumptions of Theorem~\ref{thm:charQAbelian}
throughout the present Section~\ref{ssec:pf44}. The proof follows the strategy
of \cite{MR3030068} closely, but has to overcome a fair number of technical
difficulties arising from the presence of singularities. It is fairly long and
therefore subdivided into numerous steps. The main idea is to show abundance
for $X$, and then to analyse the Iitaka fibration. Abundance is shown in Steps
4, 5 and 6, with an inductive argument using repeated covers, fibrations, and
restrictions to fibres.

\begin{itemize}
\item Steps 1--3 set the stage, prove minimality of $X$ and put limits on its
possible numerical dimension and Kodaira dimension.

\item Step 4 fibres $X$ using a Shafarevich map construction. The general fibre
is called~$F$.

\item Step 5, which is the longest step of this proof, considers a suitable
cover of $\what{F}$ of~$F$ and takes a general fibre of its Albanese map. An
analysis of this fibre will then give abundance for $F$.

\item Step 6 uses abundance for $F$ to prove abundance for $X$.

\item Steps 7--8 end the proof by showing that the Iitaka fibration of a
suitable cover of $X$ is birational (and then isomorphic) to an Abelian group
scheme over a proper base.
\end{itemize}

\subsection*{Step 1: Simplification}

Since assumption and conclusion of Theorem~\ref{thm:charQAbelian} are stable
under quasi-étale covers, we may apply \cite[Th.\,1.5]{GKP13}, pass to a maximal
quasi-étale cover and assume that the following holds in addition.

\begin{asswlog}[$X$ is maximally quasi-étale]\label{asswlog:mqe}
The variety $X$ is maximally quasi-étale. In other words, the algebraic
fundamental groups $\what{\pi}_1(X_{\reg})$ and $\what{\pi}_1(X)$ agree.
\end{asswlog}

With this assumption in place, the criterion for projective flatness,
\cite[Th.\,1.6]{GKP20a}, implies that $\Omega^1_{X_{\reg}}$ is projectively flat. As
we have seen, this has a number of interesting consequences.

\begin{consequence}[Extension of projectively flat bundles]\label{cons:5-3}
The extension result for projectively flat bundles, \cite[Prop.\,3.10]{GKP20a},
allows to find a projectively flat $\bP^{n-1}$-bundle $\bP \to X$ and an isomorphism
of $X$-schemes, $\bP|_{X_{\reg}} \cong_X \bP\big(\Omega^1_{X_{\reg}} \bigr)$.
\end{consequence}

\begin{consequence}[Quotient singularities]\label{consequence:isolated}
The local description of projectively flat sheaves, \cite[Prop.\,3.11]{GKP20a},
and Proposition~\ref{prop:icq} imply that $X$ has only isolated cyclic
quotient singularities.
\end{consequence}

\begin{notation}\label{not:5-5}
We chose a projectively flat bundle $\bP \to X$ as in Consequence~\ref{cons:5-3}
and maintain this choice throughout. We write $\tau: \pi_1(X) \to \PGL(n,\bC)$ for the
representation that defines the projectively flat structure on $\bP$.
\end{notation}

If $K_X$ is numerically trivial, then the Chern class equality \eqref{eq:qcce}
reduces to \hbox{$\what{c}_2(X)\cdot [H]^{n-2} = 0$}, and we may apply
\cite[Th.\,1.2]{LT18} to conclude that $\what{X}$ is quasi-Abelian. This allows
to make the following assumption.

\begin{asswlog}[$K_X \not \equiv 0$]\label{asswlog:knnt}
The canonical class $K_X$ is not numerically trivial.
\end{asswlog}

\subsection*{Step 2: Minimality of $X$}

Next, we show that the variety $X$ is minimal.

\begin{claim}[Minimality of $X$]\label{claim:minimal}
The canonical divisor $K_X$ is nef.
\end{claim}
\begin{proof}[Proof of Claim~\ref{claim:minimal}]
Choose an integer $m$ such that $\omega^{[\otimes m]}_X$ is invertible. We argue by
contradiction and assume that there exists a contraction of a $K_X$-negative
extremal ray. Choosing a resolution of $X$, we find a sequence of morphisms
\[
\wtilde{X} \To{\pi\text{, resolution}} X \To{\phi\text{, contraction}}Y.
\]
If $y \in Y$ is any point, Takayama's result on the local fundamental groups of
$\phi$, Proposition~\ref{prop:takayama}, allows to find a neighbourhood
$U = U(y) \subseteq Y^{(\an)}$, open in the analytic topology, such that
$V := \phi^{-1}(U) \subseteq X^{(\an)}$ is connected and simply connected. Using simple
connectedness, the local description of projectively flat sheaves,
\cite[Prop.\,3.11]{GKP20a}, provides us with a reflexive, rank-one coherent
analytic sheaf $\sL$ on $V$ and an isomorphism $\Omega^{[1]}_V \cong \sL^{\oplus n}$. We are
done once we show that the first Chern class of the invertible sheaf
$\sL^{[\otimes n\cdot m]}$,
\[
c_1(\sL^{[\otimes n\cdot m]}) \in H^1\bigl(V,\, \Omega^1_V \bigr),
\]
intersects all curves trivially, as this will contradict negativity of the
contracted extremal ray. To this end, write $\wtilde{V} := \pi^{-1}(V)$ and
consider the following commutative diagram,
\begin{equation}\label{eq:dDruel}
\begin{tikzcd}
\Pic(V) \ar[r, "\ts c_1"] \ar[d, "\ts \pi^*"'] & H^1\bigl(V,\, \Omega^1_V \bigr) \ar[r] \ar[d, "\ts \diff\pi"] & H^1\bigl(V,\, \Omega^{[1]}_V \bigr) \ar[r, "\ts \alpha"] \ar[d, "\ts\beta"] & H^1\bigl(V_{\reg},\, \Omega^1_{V_{\reg}} \bigr) \\
\Pic(\wtilde{V}) \ar[r, "\ts c_1"'] & H^1\bigl(\wtilde{V},\, \Omega^1_{\wtilde{V}} \bigr) \ar[r, equal] & H^1\bigl(\wtilde{V},\, \Omega^1_{\wtilde{V}} \bigr)
\end{tikzcd}
\end{equation}
where $\beta$ is induced by the pull-back map for reflexive differentials, and
where $c_1$ maps a bundle with transition functions $f_{ij}$ to
$\{d \log f_{ij}\}_{i,j}$. We have seen in Proposition~\ref{prop:druel} that
the image of $c_1(\sL^{[\otimes n\cdot m]})$ in $H^1\bigl(V_{\reg},\, \Omega^1_{V_{\reg}} \bigr)$
vanishes. As $\Omega^{[1]}_X$ is reflexive, the morphism $\alpha$ is injective by
\cite[Th.\,1.14]{MR0287033} and \cite[Cor.\,II.3.15]{BS76}; see also
\cite[Th.\,II.3.6]{BS76}. By commutativity of \eqref{eq:dDruel}, the Chern
class $c_1(\pi^* \sL^{[\otimes n\cdot m]})$ of the pullback via $\pi$ vanishes. The claim
follows from the projection formula for intersection numbers.
\end{proof}

\subsection*{Step 3: Kodaira dimension}

We aim to show that the minimal variety $X$ is quasi-Abelian, which implies in
particular that $\kappa(X)=0$. As a first step, we show that $X$ is at least not of
general type.

\begin{claim}\label{claim:ngen}
The variety $X$ is not of general type.
\end{claim}
\begin{proof}[Proof of Claim~\ref{claim:ngen}]
Suppose to the contrary that $X$ \emph{is} of general type. Using that $X$ is
also minimal, recall from \cite[Th.\,1.1]{GKPT19b} that the Chern classes of
$X$ satisfy a $\bQ$-Miyaoka-Yau inequality,
\begin{equation}\label{eq:ace1}
\frac{n}{2(n+1)}\cdot [K_X]^n = %
\frac{n}{2(n+1)}\cdot \what{c}_1\bigl(\Omega^{[1]}_X \bigr)^2\cdot [K_X]^{n-2} %
\leq \what{c}_2\bigl(\Omega^{[1]}_X \bigr)\cdot [K_X]^{n-2}.
\end{equation}
On the other hand, we assume that Equation~\eqref{eq:qcce} holds. An
elementary computation, using \cite[(1.14) on p.\,34]{Kob87} and
\cite[\S 3.8]{GKPT19b} shows that Equation~\eqref{eq:qcce} is equivalent to
the assertion that the semistable sheaf $\sEnd(\Omega^{[1]}_X)$ has vanishing
$\bQ$-Chern classes with respect to $H$, in the sense of
\cite[Defn.~6.1]{GKPT19}. But then we have seen in \cite[Fact and
Defn.~6.5]{GKPT19} that $\sEnd(\Omega^{[1]}_X)$ has vanishing $\bQ$-Chern classes
with respect any ample bundle, and hence also with respect to any nef bundle.
We obtain that
\begin{equation}\label{eq:ace2}
\frac{n-1}{2n}\cdot [K_X]^n %
= \frac{n-1}{2n}\cdot \what{c}_1 \bigl(\Omega^{[1]}_X \bigr)^2\cdot [K_X]^{n-2} %
= \what{c}_2 \bigl(\Omega^{[1]}_X \bigr)\cdot [K_X]^{n-2}.
\end{equation}
But we know that $K_X$ is big and nef, so that $[K_X]^n > 0$. Putting
(in)equalities \eqref{eq:ace1} and \eqref{eq:ace2} together hence produces a
contradiction. \qedhere~\mbox{(Claim~\ref{claim:ngen})}
\end{proof}

\subsection*{Step 4: The Shafarevich map}

Recall Notation~\ref{not:5-5} and consider the representation
$\tau: \pi_1(X) \to \PGL(n,\bC)$ that defines the projectively flat structure on $\bP$.
Write
\[
G_X = \overline{\img(\tau)} \subseteq \PGL(n,\bC)
\]
for the algebraic Zariski closure
of its image. This is a linear algebraic group, which has finitely many
components. Applying Selberg's Lemma\footnote{see \cite{MR925989} for a brief
account of this classic result.} and passing to an appropriate étale cover of
$X$, we may hence assume without loss of generality that the following holds.

\begin{asswlog}[Connectivity and torsion-freeness]\label{asswlog:conn}
The linear algebraic group $G_X$ is connected. Writing $\Rad(G_X) \lhd G_X$
for the solvable radical of $G_X$, the image of the composed group
homomorphism to the linear algebraic group $\factor{G_X}{\Rad(G_X)}$ is
torsion-free,
\begin{equation}\label{eq:5-10-1}
\pi_1(X) \To{\tau} G_X \to \factor{G_X}{\Rad(G_X)}.
\end{equation}
\end{asswlog}
If $G_X$ reduces to a point, we have $\Omega^{[1]}_X \cong \sL^{\oplus n}$ for a suitable
Weil-divisorial sheaf~$\sL$. We may then apply Theorem~\ref{thm:druel} to
conclude that $X$ is quasi-Abelian. As a consequence, we may assume without
loss of generality that the following holds.

\begin{asswlog}[Representation has infinite image]\label{claim:5-x}
The group $G_X$ is of positive dimension.
\end{asswlog}

We now follow the strategy of \cite{MR3030068} and use
Assumption~\ref{claim:5-x} to consider the Shafarevich map for the solvable
radical of $G_X$. We refer the reader to \cite[\S 3]{MR1341589} for more on
Shafarevich maps.

\begin{construction}
Consider the composed morphism \eqref{eq:5-10-1} of
Assumption~\ref{asswlog:conn}, let $K \lhd \pi_1(X)$ be its kernel and consider
an associated $K$-Shafarevich map, $\sh^K(X): X \dasharrow \Sh^K(X)$, where
$\sh^K(X)$ is dominant and $\Sh^K(X)$ is a smooth, projective variety.
\end{construction}

The following two properties of $\sh^K(X)$ will be most relevant in the sequel.

\begin{fact}[\protect{Shafarevich maps are almost holomorphic, \cite[Th.\,3.6]{MR1341589}}]
The rational map $\sh^K(X)$ is almost holomorphic. In other words, there
exists a Zariski open set $X^\circ \subseteq X$ such that $\sh^K(X)|_{X^\circ}$ is well-defined
and proper. The fibres of $\sh^K(X)|_{X^\circ}$ are connected. \qed
\end{fact}

\begin{fact}[\protect{Shafarevich maps and fundamental groups, \cite[Th.\,3.6]{MR1341589}}]\label{fact:Sh2a}
Let $x \in X^\circ$ be a very general point and let $Z \subset X$ be a subvariety through
$x$, with normalisation $\eta : \wtilde{Z} \to Z$. Then, the rational map
$\sh^K(X)$ maps $Z$ to a point if and only if the composed morphism
\[
\pi_1(\wtilde{Z}) \To{\eta_*} \pi_1(X) \To{\tau} G_X \to
\factor{G_X}{\Rad(G_X)}
\]
has finite image. \qed
\end{fact}

\begin{claim}[The base of the Shafarevich map]\label{claim:6-30}
The (smooth) variety $\Sh^K(X)$ is of general type.
\end{claim}
\begin{proof}[Proof of Claim~\ref{claim:6-30}]
Choose a resolution $\pi : \wtilde{X} \to X$, recall from
\cite[Th.\,1.1]{Takayama2003} that the natural push-forward map
$\pi_* : \pi_1(\wtilde{X}) \to \pi_1(X)$ is isomorphic and observe that the composed
map $\sh^K(X) \circ \pi : \wtilde{X} \dasharrow \Sh^K(X)$ is therefore a Shafarevich
map for $K \lhd \pi_1(\wtilde{X})$. Using the assumption that the image of the
composed morphism \eqref{eq:5-10-1} is torsion-free, the claim then follows
from \cite[Th.\,1]{MR3314830}. \qedhere~\mbox{(Claim~\ref{claim:6-30})}
\end{proof}

Claims~\ref{claim:ngen} and \ref{claim:6-30} together imply in particular that
the fibres of $\sh^K(X)$ are positive-dimensional. These will be investigated
next.

\subsection*{Step 5: Fibres of the Shafarevich map}

Throughout the present step, choose a general fibre $F \subset X$ of $\sh^K(X)$.
There are two cases to consider: either $\Sh^K(X)$ is a point and $F = X$ is
potentially singular, or $F$ is a proper subvariety of $X$, and then it avoids
the (finitely many!) singularities of $X$,
cf.~Consequence~\ref{consequence:isolated}, and must hence be smooth. Either
way, we will show in this step that $F$ is quasi-Abelian. The proof is somewhat
long and therefore divided into sub-steps.

\begin{notation}
Given an étale cover $\gamma :\what{F} \to F$, consider the composed group morphism
\begin{equation}\label{eq:sgsdfg}
\pi_1\bigl(\what{F}\,\bigr) \to\pi_1(F) \to \pi_1(X) \To{\tau}
\PGL(n,\bC)
\end{equation}
and let $G_\gamma \subseteq \PGL(n,\bC)$ be the Zariski closure of its image.
\end{notation}

\begin{rem}\label{rem:connected=solvable}
Given an étale cover $\gamma :\what{F} \to F$, the group $G_\gamma$ is an algebraic
subgroup of $\PGL(n,\bC)$ with finitely many connected components.
Fact~\vref{fact:Sh2a} implies that the maximal connected subgroup of $G_\gamma$ is
solvable.
\end{rem}

\subsection*{Step 5-1: Irregularity}

To begin our analysis of the fibres of the Shafarevich map, we show that the
fibres have positive irregularity, at least once we pass to a suitable étale
cover. This will later allow to analyse the fibres by means of their Albanese
map.

\begin{claim}\label{claim:uze}
There exists an étale cover $\gamma : \what{F} \to F$ such that
$h^0 \bigl(\what{F},\, \Omega^{[1]}_{\what{F}}\bigr) > 0$.
\end{claim}
\begin{proof}[Proof of Claim~\ref{claim:uze}]
Passing to a first étale cover of $F$, we may assume without loss of
generality that $G_\gamma$ is connected. If the composed map \eqref{eq:sgsdfg} has
infinite image, then by Remark~\ref{rem:connected=solvable} this defines a
solvable representation of $\pi_1\bigl(\what{F}\,\bigr)$ with infinite image.
It follows that $H^1 \bigl(\what{F},\, \bZ \bigr)$ has positive rank and
consequently, that $\dim_\bC H^1\bigl(\what{F},\, \bC \bigr)>0$. The description
of the natural Hodge structure on this space, \cite[Th.\,1]{Schwald16}, then
implies that $h^0 \bigl(\what{F},\, \Omega^{[1]}_{\what{F}}\bigr) > 0$, as desired.

It therefore remains to consider the case where the composed map
\eqref{eq:sgsdfg} has finite image. By Assumption~\ref{claim:5-x}, this
implies that $F$ is a proper subvariety of $X$, entirely contained in
$X_{\reg}$. We may then choose $\gamma$ such that the composed map
\eqref{eq:sgsdfg} is trivial, implying as before that the pull-back of the
projectively flat sheaf $\Omega^{[1]}_X$ is a direct sum of the form,
$\gamma^* \Omega^1_X \cong \sL^{\oplus n}$, for some $\sL \in \Pic(\what{F})$. The pull-back of the
conormal bundle sequence, which is non-trivial for reasons of dimension, will
then read as follows,
\[
0 \to \underbrace{\gamma^* N^*_{F/X}}_{\text{trivial bundle}} \to
\underbrace{\gamma^* \Omega^1_X}_{\cong \sL^{\oplus n}} \to \Omega^1_{\what{F}} \to 0.
\]
Hence $H^0\bigl(\what{F},\, \sL\bigr) \ne \{0\}$, and consequently
$H^0\bigl(\what{F},\, \Omega^1_{\what{F}}\big) \ne \{0\}$.
\qedhere~\mbox{(Claim~\ref{claim:uze})}
\end{proof}

\begin{notation}\label{not:50-20}
For the remainder of this proof, fix one particular Galois cover
\hbox{$\gamma : \what{F} \to F$} such that $G_\gamma$ is connected, and where
$h^0 \bigl(\what{F},\, \Omega^{[1]}_{\what{F}}\bigr) > 0$. Recalling
Remark~\ref{rem:connected=solvable}, we may apply
Proposition~\ref{prop:pfbisr} and using its notation we write
\begin{equation}\label{eq:OmegaFA0-1a}
\gamma^* \Omega^{[1]}_X \cong \sF \otimes \sA,
\end{equation}
where $\sA$ is Weil divisorial with numerical class $[\sA] \equiv [\gamma^* K_X]$, and
where $\sF$ is locally free and admits a filtration
\begin{equation}\label{eq:5-20-2}
0 = \sF_0 \subsetneq \sF_1 \subsetneq \cdots \subsetneq \sF_n = \sF
\end{equation}
such that all sheaves $\sF_i$ are locally free and flat with $\rank (\sF_i)= i$
and such that the sheaf $\sF_1$ is trivial, $\sF_1 \cong \sO_X$. In particular, there
exists an inclusion
\begin{equation}\label{eq:gjfgjh}
\sA = \sA \otimes \sF_1 \hto \gamma^* \Omega^{[1]}_X.
\end{equation}
The pull-back of the normal bundle sequence (which is trivial if $F=X$) reads
\begin{equation}\label{eq:iopz}
0 \to \underbrace{\gamma^* N^*_{F/X}}_{\text{trivial bundle}} \to
\underbrace{\gamma^* \Omega^{[1]}_X}_{= \sF \otimes \sA} \to \Omega^{[1]}_{\what{F}} \to 0.
\end{equation}
\end{notation}

\subsection*{Step 5-2: Numerical dimension and Kodaira dimension}

Claim~\ref{claim:uze} has fairly direct consequences for the Kodaira dimension
of the fibres, which will turn out to be either zero or one, the second case
disappearing eventually.

\begin{claim}[Bounding the Kodaira dimension from below]\label{claim:Kodairadim2}
Maintaining Notation~\ref{not:50-20}, if $\pi : \wtilde{F} \to \what{F}$ is any
resolution of singularities, then the Kodaira dimension of $\wtilde{F}$ is
non-negative. In particular, it follows that $\kappa(\what{F}) \geq 0$.
\end{claim}
\begin{proof}[Proof of Claim~\ref{claim:Kodairadim2}]
The assumption that $h^0 \bigl(\what{F},\, \Omega^{[1]}_{\what{F}}\bigr) > 0$,
together with the Isomorphism~\eqref{eq:OmegaFA0-1a} and the
Filtration~\eqref{eq:5-20-2} allows to find an index $i$ such that
\begin{equation}\label{eq:ugh}
h^0\bigl(\what{F},\, \sA \otimes \factor{\sF_i}{\sF_{i-1}} \bigr) \ne 0.
\end{equation}
Writing $\sL_i := \factor{\sF_i}{\sF_{i-1}}$ and recalling that $\sL_i$ is invertible
and flat, hence numerically trivial by Lemma~\ref{lem:c1offlat},
Equation~\eqref{eq:ugh} implies in particular that the sheaf
\[
\sA^{[n]} \otimes \sL^{\otimes n}_i = \underbrace{\sA^{[n]} \otimes \det \sF\vphantom{\sL^{\otimes n}_i}}_{=
\gamma^*\omega_X = \omega_{\what{F}}} \quad \otimes \quad \underbrace{(\det \sF)^* \otimes \sL^{\otimes
n}_i}_{\mathclap{=: \sL' \text{, num.\ trivial by Rem.~\ref{rem:pfbisr}}}}
\]
also has a non-trivial section. In this setting, the extension theorem for
differential forms, \cite[Th.\,4.3]{GKKP11} yields a non-trivial morphism
$\pi^{[*]} \omega_{\what{F}} \hto \omega_{\wtilde{F}}$ and hence a non-trivial section of
$\omega_{\wtilde{F}} \otimes \pi^*\sL'$, where again $\pi^*\sL' \in \Pic(\wtilde{F})$ is
numerically trivial. The claim thus follows from the ``numerical character of
the effectivity of adjoint line bundles'', \cite[Th.\,0.1]{CKP12}, applied to
the smooth space $\wtilde{F}$.
\qedhere~\mbox{(Claim~\ref{claim:Kodairadim2})}
\end{proof}

\begin{claim}[Bounding the numerical dimension from above]\label{claim:numericaldimensionbound-x}
The numerical dimen\-sion of $\what{F}$ satisfies $\nu(\what{F}) \leq 1$.
\end{claim}
\begin{proof}[Proof of Claim~\ref{claim:numericaldimensionbound-x}]
To begin, recall from \eqref{eq:OmegaFA0-1a} that the (nef!) numerical classes
$n\cdot [\sA]$ and $[\gamma^* K_X] = [K_{\what{F}}]$ agree. It follows that
$\nu(\what{F}) = \nu(\sA)$, and so it suffices to consider $\nu(\sA)$ and to show that
$\nu(\sA) \leq 1$. In other words, given any very ample divisor $H$ on $\what{F}$,
we need to show that
\begin{equation}\label{eq:xxl}
[\sA]^2\cdot [H]^{m-2} = 0, \quad \text{where } m := \dim \what{F}.
\end{equation}

Recall from \eqref{eq:gjfgjh} that there exists an embedding
$\sA \hto \gamma^* \Omega^{[1]}_X$. Combined with Sequence~\eqref{eq:iopz}, we obtain a
sheaf morphism
\begin{equation}\label{eq:jhk}
\sA \to \Omega^{[1]}_{\what{F}}.
\end{equation}
If this morphism is constantly zero, then $\sA$ maps into the trivial sheaf
$\gamma^* N^*_{F/X}$, the nef sheaf $\sA$ must thus be trivial, and the claim is
shown. We will therefore continue this proof under the assumption that
\eqref{eq:jhk} is an embedding.

Next, choose a general tuple $(H_1, \dots, H_{m-2}) \in |H| \times \cdots \times |H|$ and consider
the associated complete intersection surface $S := H_1 \cap\cdots\cap H_{m-2}$. Then,
$S$ has klt singularities\footnote{The surface $S$ will in fact almost always be smooth, except perhaps when $X\!=\!F$ and \hbox{$\dim X\!=\!2$}.} and $\sA|_S$ is
reflexive. The kernel of the restriction morphism for reflexive differentials
is
\[
N_{S/\what{F}}^* = \ker \bigl(\Omega^{[1]}_{\what{F}}\bigr|_S \ra \Omega^{[1]}_S \bigr)
\cong \sO_{\what{F}}(-H)^{\oplus(m-2)},
\]
so in particular anti-ample. But since the nef sheaf $\sA$ never maps to an
anti-ample, we find that the composed sheaf morphism
\[
\sA|_S \to \Omega^{[1]}_{\what{F}} \Bigr|_S \to \Omega_S^{[1]}
\]
cannot vanish. The Bogomolov-Sommese vanishing theorem for the potentially
singular surface $S$, \cite[Th.\,8.3]{GKK08}, implies that $\sA|_S$ is not big.
Since it is nef, we conclude that
\[
0 = [\sA|_S]^2 = [\sA]^2\cdot [H]^{m-2},
\]
as desired. \qedhere~\mbox{(Claim~\ref{claim:numericaldimensionbound-x})}
\end{proof}

\begin{claim}[Kodaira dimension and numerical dimension of $F$]\label{claim:kF}
The Kodaira-dimension and the numerical dimension of $F$ satisfy the
inequalities
\begin{equation}\label{eq:u56}
0 \leq \kappa(F) \leq \nu(F) \leq 1.
\end{equation}
\end{claim}
\begin{proof}[Proof of Claim~\ref{claim:kF}]
Given that $\gamma : \what{F} \to F$ is étale,
Claim~\ref{claim:numericaldimensionbound-x} immediately implies that
$\nu(F) \leq 1$. Since $K_F = K_X|_F$ is nef, \cite[Prop.\,2.2]{MR782236} gives the
two rightmost inequalities in \eqref{eq:u56}, and so it remains to show that
$\kappa(F) \geq 0$. But we already know from Claim~\ref{claim:Kodairadim2} that
$\kappa(\what{F}) \geq 0$, so that there exists a number $p \in \bN$ and a non-trivial
section $\sigma \in H^0\bigl(\what{F},\, \omega^{\otimes p}_{\what{F}}\bigr)$. But then
\[
\bigotimes_{g \in \Gal(\gamma)} g^*\sigma \in H^0\bigl(\what{F},\, \omega^{\otimes p\cdot \#\Gal(\gamma)}_{\what{F}}\bigr)
\]
is a non-trivial Galois-invariant pluri-form on $\what{F}$, which hence
descends to a non-trivial pluri-form on $F$.
\qedhere~\mbox{(Claim~\ref{claim:kF})}
\end{proof}

\subsection*{Step 5-3: Fibres of the Albanese map}

We pointed out above that the Albanese map of $F$ is not trivial, at least once
we pass to a suitable cover. Using the fact that the Albanese yields a
Shafarevich map for the commutator subgroup of the fundamental group, there is
much that we can say about its fibres, which will eventually turn out to be
abundant.

\begin{construction}\label{cons:itAlb-s}
Maintaining Notation~\ref{not:50-20}, consider the Albanese map
$\alb(\what{F}) : \what{F} \to \Alb(\what{F})$ and recall from
Claim~\ref{claim:uze} and Notation~\ref{not:50-20} that this map is
non-trivial, which is to say that $\dim \Alb(\what{F}) > 0$. Let
$F_1 \subsetneq \what{F}$ be a general fibre component of $\alb(\what{F})$. This might
be a point if $\what{F}$ is of maximal Albanese dimension. Since $\what{F}$
has at worst isolated singularities, $F_1$ is necessarily smooth. Now, we use
the fact that Albanese maps are Shafarevich maps by recalling from
\cite[§0.1.3]{MR1341589} that the push-forward morphism of Abelianised
fundamental groups,
\[
(\iota_1)^{(\ab)}_* : \pi_1(F_1 )^{(\ab)} \to \pi_1\bigl(\what{F}\, \bigr)^{(\ab)}
\]
has finite image. We can thus find an étale cover $\gamma_1 : \what{F}_1 \to F_1$
such that the composed push-forward
\[
(\iota_1\circ \gamma_1)^{(\ab)}_* : \pi_1(\what{F}_1 )^{(\ab)} \to \pi_1\bigl(\what{F}\,
\bigr)^{(\ab)}
\]
is trivial. The following diagram summarises the morphisms that we have
discussed so far,
\[
\begin{tikzcd}[column sep=2cm, row sep=1cm]
\what{F}_1 \ar[rrrr, bend left=10, "\eta"] \ar[r, two heads, "\gamma_1\text{, étale}"'] & F_1 \ar[r, hook, "\iota_1\text{, fibre inclusion}"'] & \what{F} \ar[r, two heads, "\gamma\text{, étale}"'] \ar[d, "\alb(\what{F})"] & F \ar[r, hook, "\iota\text{, fibre inclusion}"'] & X \ar[d, dashed, "\sh^K\!(X)"'] \\
&& \Alb(\what{F}) && \Sh^K\!(X)
\end{tikzcd}
\]
Observing that $\img(\eta) \subsetneq X_{\reg}$, we consider the exact sequence of
differentials,
\begin{equation}\label{eq:trztu-s}
0 \to \underbrace{\Omega^1_{\what{F}_1/X}}_{\mathclap{\text{trivial,
non-zero}}} \to \eta^* \Omega^1_X \To{\diff\eta} \Omega^1_{\what{F}_1} \to0.
\end{equation}
\end{construction}

\begin{claim}[Numerical invariants of $F_1$ and $\what{F}_1$]\label{claim:6-23a-s}
Assume the setting of Construction~\ref{cons:itAlb-s}. Then,
$0 \leq \kappa(F_1) \leq \nu(F_1) \leq 1$ and $0 \leq \kappa(\what{F}_1) \leq \nu(\what{F}_1) \leq 1$.
\end{claim}
\begin{proof}[Proof of Claim~\ref{claim:6-23a-s}]
In order to lighten notation a little bit, we will suppress $\iota_1$ in the
following discussion. Given that the morphism $\gamma_1$ is étale, it suffices to
consider the variety $\what{F}_1$ only. We will also assume that $\what{F}_1$
is not a point, or else there is little to show, cf.~Claim~\ref{claim:kF}. On
the one hand, using that $\img(\eta) \subsetneq X_{\reg}$, we obtain a description of the
canonical bundle,
\[
\omega_{\what{F}_1} \overset{\text{\eqref{eq:trztu-s}}}{\cong} \eta^*(\omega_X)
\overset{\text{\eqref{eq:OmegaFA0-1a}}}{\cong} \gamma_1^* \bigl(\sA^{[\otimes n]} \otimes \det
\sF\bigr).
\]
Observing that $\gamma_1^* \sA$ is invertible and that $\gamma_1^* \sF$ is flat, this gives
a numerical equivalence $[\omega_{\what{F}_1}] \equiv n\cdot \bigl[\gamma_1^* \sA\bigr]$. On
the other hand, \eqref{eq:gjfgjh} and \eqref{eq:trztu-s} combine to give a
morphism
\[
\alpha : \gamma_1^* \sA \to \Omega^1_{\what{F}_1}.
\]
There are two cases to consider. If the morphism $\alpha$ is not zero,
Claim~\ref{claim:6-23a-s} follows from \cite[Lem.\,3.1]{MR3030068}. Otherwise,
if the morphism $\alpha$ is zero, we obtain an embedding of the nef line bundle
$\gamma_1^* \sA$ into the trivial sheaf $\Omega^1_{\what{F}_1/X}$. It follows that
$\gamma_1^* \sA$ and hence $\omega_{\what{F}_1}$ are numerically trivial, so
$\nu(\what{F}_1) = 0$. By \cite[Th.\,8.2]{Kawamata85}, this implies
$\kappa(\what{F}_1) = 0$. Either way, Claim~\ref{claim:6-23a-s} follows.
\qedhere~\mbox{(Claim~\ref{claim:6-23a-s})}
\end{proof}

\begin{claim}[Semiampleness of $K_{\what{F}_1}$ and $K_{F_1}$]\label{claim:6-25a-s}
Assume the setting of Construction~\ref{cons:itAlb-s}. Then, the canonical
bundles $K_{\what{F}_1}$ and $K_{F_1}$ are semiample.
\end{claim}
\begin{proof}[Proof of Claim~\ref{claim:6-25a-s}]
As before, we consider the variety $\what{F}_1$ only and assume that it is not
a point. We also consider the following commutative diagram of group
morphisms,
\[
\begin{tikzcd}[column sep=2cm, row sep=1cm]
\pi_1\bigl(\what{F}_1 \bigr) \ar[r, "\ts\tau\circ \eta_*"] \ar[d, "\underset{\scriptstyle\text{trivial after Abelianisation}}{\ts(\iota_1\circ \gamma_1)_*}"'] & G_X \ar[r, hook, "\ts\text{subgroup}"] & \PGL(n,\bC) \\
\pi_1\bigl(\what{F} \bigr) \ar[r, "\ts\tau\circ (\iota\circ \gamma)_*"'] & G_\gamma \ar[u, hook, "\ts\text{subgroup}"']
\end{tikzcd}
\]
The representation $\tau\circ \eta_*$ is solvable because it factors via the solvable
group $G_\gamma$. Better still, the representation $\tau\circ \eta_*$ is in fact trivial, or
else there would be a non-trivial Abelian subrepresentation, contradicting
triviality of the Abelianised push-forward $(\iota_1\circ \gamma_1)^{(\ab)}_*$. As a
consequence, we find that $\eta^* \Omega^1_X$ is isomorphic to a direct sum of line
bundles,
\[
\exists\, \sL \in \Pic(\what{F}_1) : \eta^* \Omega^1_X \cong \sL^{\oplus n}.
\]
Recalling from Sequence~\eqref{eq:trztu-s} that $\eta^* \Omega^1_X$ contains the
trivial, non-zero subbundle $\Omega^1_{\what{F}_1/X}$, we find that $\sL$ admits a
section. But then $\eta^* \Omega^1_X$ is generically generated, and so is its quotient
$\Omega^1_{\what{F}_1}$. By \cite[Rem.~2.2]{Fuj13}, the manifold $\what{F}_1$ is
then of maximal Albanese dimension. But then \cite[Th.\,4.2]{Fuj13} implies
that the nef canonical bundle $K_{\what{F}_1}$ is semiample, as claimed.
\qedhere~\mbox{(Claim~\ref{claim:6-25a-s})}
\end{proof}

\subsection*{Step 5-4: Abundance and description of $F$}

In this step, we finally prove abundance for $F$ and describe its Iitaka
fibration. According to Claim~\ref{claim:kF}, there are only two cases,
$\kappa(F) = 0$ and $\kappa(F) = 1$, which we consider separately.

\begin{claim}\label{claim:solvK0}
If $\kappa(F)=0$, then $F$ is quasi-Abelian.
\end{claim}
\begin{proof}[Proof of Claim~\ref{claim:solvK0}]
To keep this proof readable, we consider three cases separately.
\begin{enumeraten}
\item\label{il:A} The divisor $K_F$ is numerically trivial and $F = X$.

\item\label{il:B} The divisor $K_F$ is numerically trivial and $F \subsetneq X$.

\item\label{il:C} The divisor $K_F$ is not numerically trivial.
\end{enumeraten}
Case~\ref{il:A} is easiest. If $K_X$ is numerically trivial, then
Equation~\eqref{eq:qcce} guarantees that
\[
\what{c}_2 \bigl(\Omega^{[1]}_X \bigr)\cdot [H]^{n-2} = 0.
\]
Then, $X$ is quasi-Abelian by \cite[Th.\,1.2]{LT18} and we are done. In
Case~\ref{il:B}, apply Item~\ref{il:6-5-2} of Lemma~\ref{lem:6-5} to the
inclusion $F \subset X_{\reg}$ to find again that $F$ is quasi-Abelian.

We consider Case~\ref{il:C} for the remainder of the proof, with the aim of
producing a contradiction. To begin, we consider a strong log resolution of
$\what{F}$ and the Stein factorisation of the associated Albanese morphism.
We obtain a commutative diagram as follows
\[
\begin{tikzcd}[column sep=2.5cm]
& \wtilde{F} \ar[r, "\ts a\text{, conn.~fibres}"] \ar[d, "\ts\pi\text{, strong log res.}"'] & S \ar[r, "\ts b\text{, finite}"] & \Alb(\wtilde{F}) \ar[d, "\ts\cong"] \\
F & \what{F} \ar[rr, "\ts \alb_{\what{F}}"'] \ar[l, "\ts\gamma\text{, étale}"] && \Alb(\what{F})
\end{tikzcd}
\]
By Claim~\ref{claim:uze}, $\Alb(\what{F})$ is not a point. A general fibre
$\wtilde{F}_1 \subsetneq \wtilde{F}$ of $a$ is therefore smooth.

Observe that the image $F_1 = \pi(\wtilde{F_1})$ is a connected component of a
general fibre of $\alb_{\what{F}}$, which avoids the (finitely many)
singularities of $\what{F}$. The restriction
\hbox{$\pi|_{\wtilde{F}_1} : \wtilde{F}_1 \to F_1$} is thus isomorphic, and
Claim~\ref{claim:6-25a-s} applies to show that $\wtilde{F}_1$ is a good
minimal model. More is true. Since $\kappa (F) = 0$ by assumption, we deduce from
Claim~\ref{claim:Kodairadim2} that $\kappa(\wtilde{F}) = 0$. Hence it follows from
\cite[Th.\,4.2]{Lai11} that $\wtilde{F}$ has a good minimal model
$\what{F}_{\min}$, which has terminal singularities and numerically trivial
canonical class.

We use the existence of $\what{F}_{\min}$ to describe $K_{\what{F}}$ in more
detail. To this end, resolve the singularities of the rational map
$\what{F}_{\min} \dasharrow \what{F}$ and obtain morphisms as follows,
\[
\begin{tikzcd}[column sep=1.5cm]
\what{F}_{\min} \ar[d, "\ts\alb_{\what{F}_{\min}}"'] & Y \ar[r, "\ts\pi_1", "\ts\text{birational}"'] \ar[l, "\ts\pi_2"', "\ts\text{birational}"] \ar[d, "\ts\alb_{Y}"] & \what{F} \ar[d, "\ts\alb_{\what{F}}"] \\
\Alb(\what{F}_{\min}) \ar[r, equal] & \Alb(Y) \ar[r, equal] & \Alb(\what{F})
\end{tikzcd}
\]
The fact that $\what{F}_{\min}$ has terminal singularities implies that
$K_Y \equiv \what{E}$, where $\what E \in \bQ\Div(Y)$ is effective with
$\pi_2$-exceptional support. Recalling that every fibre of $\pi_2$ is rationally
connected by \cite[Cor.\,1.5]{HMcK07}, we find that every component of
$\supp \what{E}$ is uniruled. The canonical class $K_{\what{F}}$ is then
numerically equivalent to $E := (\pi_1)_*\what{E}$. The $\bQ$-divisor $E$ is
effective, nef and not trivial (because $K_{\what{F}}$ is not numerically
trivial), and its components are again uniruled. But the Albanese map
$\alb_{\what{F}}$ contracts all rational curves in $\supp E$! We claim that
the image set $\alb_{\what{F}}(\supp E)$ cannot contain any isolated points.
\begin{itemize}
\item If $\dim \alb_{\what{F}}(\what{F}) \geq 2$, this follows from nefness of
$E \equiv K_{\what{F}} = \gamma^* K_X$ and from the Negativity
Lemma~\vref{lem:negativity1}.

\item If $\dim \alb_{\what{F}}(\what{F}) = 1$, this follows from nefness of
$E \equiv K_{\what{F}} = \gamma^* K_X$, from the assumption that
$0 = \kappa(F) = \kappa(\what{F})$ and from the Negativity
Lemma~\vref{lem:negativity2}.
\end{itemize}
Either way, we find that $E$ contains rational curves that stay away from the
finitely many singular points of $\what{F}$. Since $K_X$ is nef, this
contradicts Item~\ref{il:6-5-1} of Lemma~\ref{lem:6-5}.
\qedhere~\mbox{(Claim~\ref{claim:solvK0})}
\end{proof}

\begin{claim}[Description of $F$ if $\kappa(F)=1$]\label{claim:solvK1}
If $\kappa(F)=1$, then $K_F$ is semiample and the general fibre of the Iitaka
fibration is quasi-Abelian.
\end{claim}
\begin{proof}[Proof of Claim~\ref{claim:solvK1}]
If $\kappa(F)=1$, then $\nu(F) = \kappa(F)$ by Claim~\ref{claim:kF}. It follows that
$K_F$ is semiample, \cite[Cor.\,6-1-13]{KMM87}. As $F$ has isolated
singularities at worst, the general fibre $A$ of the associated Iitaka
fibration $F \to B$ is smooth, contained in $X_{\reg}$, and has numerical
trivial canonical class. As before, use Item~\ref{il:6-5-2} of
Lemma~\ref{lem:6-5} to find that $A$ is quasi-Abelian.
\qedhere~\mbox{(Claim~\ref{claim:solvK1})}
\end{proof}

\subsection*{Step 6: Abundance for $X$, description of the Iitaka fibration}

We have seen in the previous step that the general fibres of the Shafarevich map
are abundant. As we will see now, this implies Abundance for $X$.

\begin{claim}[Abundance for $X$]\label{claim:AbdX}
The canonical bundle $K_X$ is semiample.
\end{claim}
\begin{proof}[Proof of Claim~\ref{claim:AbdX}]
If $\Sh^K(X)$ is a point, then $F = X$ and we have seen in
Claim~\ref{claim:kF} that $0 \leq \kappa(F) = \kappa(X) \leq 1$. We know that abundance holds
because the two cases have been described in Claims~\ref{claim:solvK0} and
\ref{claim:solvK1}, respectively. If $\Sh^K(X)$ is positive-dimensional, then
we can apply Proposition~\ref{prop:abundance} to the Shafarevich map.
\begin{itemize}
\item Assumption~\ref{il:AS1} is satisfied by Claim~\ref{claim:6-30}

\item Assumption~\ref{il:AS2} is satisfied as nefness of $K_X$ has been shown
in Claim~\ref{claim:minimal}.

\item Assumption~\ref{il:AS3} holds by Consequence~\ref{consequence:isolated}.

\item Assumption~\ref{il:AS4} has been shown in Claims~\ref{claim:solvK0} and
\ref{claim:solvK1}, respectively.
\end{itemize}
The claim thus follows. \qedhere~\mbox{(Claim~\ref{claim:AbdX})}
\end{proof}

\begin{rem}
Claim~\ref{claim:AbdX} implies in particular that
\begin{itemize}
\item the Iitaka fibration has a positive-dimensional base by
Assumption~\ref{asswlog:knnt}, and

\item the fibres of the Iitaka fibration are quasi-Abelian by
Item~\ref{il:6-5-2} of Lemma~\ref{lem:6-5}.
\end{itemize}
\end{rem}

We denote the Iitaka fibration by $\iit(X) : X \to \Iit(X)$.

\begin{rem}[Ambro’s canonical bundle formula]\label{rem:ACBF}
Ambro’s canonical bundle formula for projective klt pairs,
\cite[Th.\,4.1]{MR2134273} or \cite[Th.\,3.1]{MR2944479}, provides us with an
effective $\bQ$-divisor $\Delta$ on $\Iit(X)$ that makes the pair
$\bigl(\Iit(X), \Delta \bigr)$ klt.
\end{rem}

\subsection*{Step 7: Abelian group schemes}

Following the ideas of Jahnke-Radloff, \cite{MR3030068}, we~use Kollár's
characterisation of étale quotients of Abelian group schemes to show that the
Iitaka fibration of a suitable étale cover of $X$ is birational to an Abelian
group scheme.

\begin{claim}\label{claim:glfg}
There exists an étale cover $\what{X} \to X$ whose Iitaka fibration is
birational to an Abelian group scheme over a smooth projective base admitting
a level three structure. More precisely, there exists a commutative diagram
\begin{equation}\label{eq:glfg}
\begin{tikzcd}[column sep=2.5cm]
A \ar[r, dashed, "\ts\phi\text{, birational}"] \ar[d, "\ts\alpha"] & \what{X} \ar[r, "\ts\text{étale}"] \ar[d, "\ts\iit(\what{X})"] & X \\
S \ar[r, "\ts\psi\text{, birational}"'] \ar[u, bend left, "\ts\sigma"] & \Iit(\what{X})
\end{tikzcd}
\end{equation}
where $S$ is smooth and projective, and $\alpha$ is a family of polarised Abelian
varieties admitting a level three structure.
\end{claim}
\begin{proof}[Proof of Claim~\ref{claim:glfg}]
We construct from right to left a diagram as follows,
\begin{equation}\label{eq:K12}
\begin{tikzcd}[column sep=2.5cm]
A \ar[d, "\ts\alpha"'] \ar[r, "\ts\text{étale}"] & \wcheck{A} \ar[d, "\ts\wcheck{\alpha}"'] \ar[r, dashed, "\ts\text{birational}"] & \wcheck{X} \ar[d, "\ts\wcheck{a}"'] \ar[r, "\ts\varepsilon\text{, étale}"] & X \ar[d, "\ts\iit(X)"] \\
S \ar[r, "\ts\varepsilon_S \text{, étale}"'] & \wcheck{S} \ar[r, "\ts\rho\text{,
birational}"'] & \wcheck{Y} \ar[r, "\ts\beta\text{, finite}"'] & \Iit(X)
\end{tikzcd}
\end{equation}
where $\wcheck{a}$ has connected fibres, where~$\alpha$ and $\wcheck{\alpha}$ are
Abelian group schemes and where~$\alpha$ has the additional structure of a family
of Abelian varieties with level three structure. To begin, we claim that the
variety $X$ has generically large fundamental group along the general fibre
$X_y$ of the Iitaka fibration. In case where $X$ is smooth, this claim has
been shown in \cite[Prop.\,5.1]{MR3030068}. We leave it to the reader to check
that in our case, where $X$ has only finitely many singularities, the proof of
\cite[Prop.\,5.1]{MR3030068} still applies verbatim. Given that the fibres of
the Iitaka fibration are quasi-Abelian, Proposition~\ref{prop:K63} will
therefore apply to yield the right and middle square of Diagram~\eqref{eq:K12}
--- strictly speaking, Proposition~\ref{prop:K63} gives $\rho$ only as a rational
map, but we can always blow up to make it a morphism, and then pull back
$\wcheck{A}$. Once this is done, we can find a suitable étale cover
$S \to \wcheck{S}$ such that the pull-back $A := \wcheck{A} \times_{\wcheck{S}} S$ is
an Abelian group scheme with level three structure, and a result of
Grothendieck allows to equip $A/S$ with a polarisation,
cf.~\cite[Th.\,XI.1.4]{MR0260758}. For later reference, we note that the
following property holds by construction.
\begin{enumeraten}
\item\label{il:birat} If $s \in S$ is any general point, with image
$y \in \Iit(X)$, then the induced map between fibres, $A_s \dasharrow X_y$
is birational.
\end{enumeraten}

We still need to construct $\what{X}$. We begin its construction with the
observation that $K^{[m]}_X = \iit(X)^* (H)$ for a suitable number $m \in \bN^+$
and a suitable ample $H \in \Div\bigl(\Iit(X)\bigr)$. But then,
\begin{equation}\label{eq:hjfg}
K^{[m]}_{\wcheck{X}} = \varepsilon^* K^{[m]}_X = \varepsilon^* \iit(X)^* (H) = \wcheck{a}^{\,*} (\beta^* H),
\end{equation}
where $\beta^* H$ is ample on $\what{Y}$; this shows that $\wcheck{a}$ is the
Iitaka fibration for $\wcheck{X}$ and allows to apply Ambro's canonical bundle
formula, Remark~\ref{rem:ACBF}, which shows that $\wcheck{Y}$ is the
underlying space of a klt pair. We are interested in this, because
\cite[Th.\,1.1]{Takayama2003} then asserts that
$\pi_1\bigl(\wcheck{S}\,\bigr) = \pi_1\bigl(\wcheck{Y}\,\bigr)$. As a
consequence, we find that the Stein factorisation of the morphism
$S \to \wcheck{Y}$ factors via an étale cover of $\wcheck{Y}$ and therefore
gives a diagram,
\[
\begin{tikzcd}[column sep=2.5cm]
A \ar[r, "\ts\phi"', dashed] \ar[rr, dashed, bend left=15] \ar[d, "\ts\alpha"'] & \what{X} \ar[r, "\ts\text{étale}"'] \ar[d] & \wcheck{X} \ar[d, "\ts\wcheck{a}"] \\
S \ar[r, "\ts\psi \text{, birational}"] \ar[rr, "\ts\rho\circ \varepsilon_S"', bend right=15] & \what{Y} \ar[r, "\ts\varepsilon_Y\text{, étale}"] & \wcheck{Y}
\end{tikzcd}
\]
where $\what{X} := \wcheck{X} \times_{\wcheck{Y}} \what{Y}$ and where the map $\phi$
is induced by the universal property of the fibre product. Observe, exactly
as in \eqref{eq:hjfg}, that the natural projection map $\what{X} \to \what{Y}$
is the Iitaka fibration of $\what{X}$,
$\iit(\what{X}) : \what{X} \to \Iit(\what{X}) = \what{Y}$.
Assertion~\ref{il:birat} implies that $\phi$ is birational.
\qedhere~\mbox{(Claim~\ref{claim:glfg})}
\end{proof}

\begin{claim}\label{claim:glfg2}
There exists an étale cover $\what{X} \to X$ whose Iitaka fibration is
birational to an Abelian group scheme over $\Iit(\what{X})$ with level three
structure. More precisely, there exists a commutative diagram
\begin{equation}\label{eq:glfg2}
\begin{tikzcd}[column sep=2.5cm]
A \ar[r, dashed, "\ts\phi\text{, birational}"] \ar[d, "\ts\alpha"] & \what{X} \ar[d, "\ts\iit(\what{X})"] \ar[r, "\ts\text{étale}" ]& X \\
\Iit(\what{X}) \ar[r, equal] \ar[u, bend left, "\ts\sigma"] & \Iit(\what{X})
\end{tikzcd}
\end{equation}
where $\alpha$ is a family of polarised Abelian varieties with level three
structure.
\end{claim}
\begin{proof}[Proof of Claim~\ref{claim:glfg2}]
Consider Diagram~\eqref{eq:glfg} and recall that polarised Abelian varieties
with level three structure admit a fine moduli space $\sA_3$ with universal
family $\sU_3 \to\sA_3$, and that moreover $\sA_3$ does not contain any rational
curve \cite[Lem.\,5.9.3]{Kollar93}. On the other hand, while observing
Remark~\ref{rem:ACBF} recall from \cite[Cor.\,1.5]{HMcK07} that the fibres of
the morphism $\psi$ are rationally chain connected. It follows that the moduli
map $S \to \sA_3$ factors via $\psi$ to give a morphism $\Iit(\what{X}) \to \sA_3$. To~conclude, set $A := \sU_3 \times_{\sA_3} \Iit(\what{X})$.
\qedhere~\mbox{(Claim~\ref{claim:glfg2})}
\end{proof}

\subsection*{Step 8: End of proof}

We end the proof by showing that $\what{X}$ itself is an Abelian group scheme
over a smooth base with ample canonical bundle.

\begin{claim}\label{claim:6-38}
In the setting of Claim~\ref{claim:glfg2}, the rational map
$\delta := \phi^{-1} : \what{X} \dasharrow A$ is a morphism.
\end{claim}
\begin{proof}[Proof of Claim~\ref{claim:6-38}]
Given that $\what{X}$ is klt and that the fibres of $\alpha$ do not contain any
rational curves, the claim follows from \cite[Cor.\,1.6]{HMcK07}.
\qedhere~\mbox{(Claim~\ref{claim:6-38})}
\end{proof}

\begin{claim}\label{claim:6-39}
The base variety $\Iit(\what{X})$ is smooth. In particular, $A$ is smooth.
\end{claim}
\begin{proof}[Proof of Claim~\ref{claim:6-39}]
Assume not. Let $s \in \Iit(\what{X})_{\sing}$ be any singular point, and let
$a \in \alpha^{-1}(s)$ be a general point of the fibre. Note that the morphism $\alpha$
is smooth, so~$a$ will be a singular point of $A$. On the other hand, note
that the general choice of~$a$ guarantees that $\delta^{-1}(a)$ does not contain
any of the (finitely many) singular points of~$X$. It follows that $\delta$
resolves the singularity at $a$, and that the fibre $\delta^{-1}(a)$ is thus
necessarily positive-dimensional.

There is more that we can say. Recalling from Remark~\ref{rem:ACBF} that
$\bigl(\Iit(X), \Delta \bigr)$ is klt for a certain $\bQ$-divisor $\Delta$ it follows
from smoothness of $\alpha$ that $(A, \alpha^*\Delta)$ is klt as well. But then
$\delta^{-1}(a) \subsetneq \iit(\what{X})^{-1}(s) \cap \what{X}_{\reg}$ is covered by rational
curves, violating Item~\ref{il:6-5-1} of Lemma~\ref{lem:6-5} from above.
\qedhere~\mbox{(Claim~\ref{claim:6-39})}
\end{proof}

\begin{claim}\label{claim:6-40}
The morphism $\delta$ is an isomorphism and $\what{X}$ is therefore smooth.
\end{claim}
\begin{proof}[Proof of Claim~\ref{claim:6-40}]
If $a \in A$ is any (smooth!) point where $\delta$ is not isomorphic, then
$\delta^{-1}(a)$ is positive-dimensional, and every section of $K_{\what{X}}$ will
necessarily vanish along $\delta^{-1}(a)$. That contradicts the semiampleness of
$K_{\what{X}}$, Claim~\ref{claim:AbdX}.
\qedhere~\mbox{(Claim~\ref{claim:6-40})}
\end{proof}

In the current, now nonsingular setup, the following has already been observed
in \cite[Proof of Lem.\,5.1]{MR3030068}.

\begin{claim}\label{claim:6-41}
The canonical bundle of $\Iit(\what{X})$ is ample.
\end{claim}
\begin{proof}[Proof of Claim~\ref{claim:6-41}]
The smooth variety $\Iit(\what{X})$ is of general type, while the existence of
$\sigma$ and minimality of $X$ and hence $\what X$ together imply that its
canonical bundle is nef, hence abundant. Fibres of the Iitaka fibration for
$\Iit(\what{X})$ are covered by rational curves, which do not exist by the
existence of $\sigma$ and the fact that $\what X$ (as an étale cover of $X$) does
not contain rational curves by Item~\ref{il:6-5-1} of Lemma~\ref{lem:6-5}.
\end{proof}

Now that we know that $\what{X}$ is an Abelian group scheme over a smooth
projective base with ample canonical bundle, the final step of the argument of
Jahnke-Radloff, \cite[Th.\,6.1]{MR3030068}, applies to conclude also our proof.
\qed~\mbox{(Theorem~\ref{thm:charQAbelian})}

\section{Proof of Theorem~\ref*{thm:6-1}}

\subsection*{Step 1: Preparations}

In the setup of Theorem~\ref{thm:6-1}, we set
\[
\sF := \bigl(\Sym^n\Omega^1_X \otimes \sO_X(-K_X)\bigr)^{**}.
\]
Let $H$ be a very ample divisor on $X$ and let
$(D_1, \dots, D_{n-1}) \in |H|^{\times (n-1)}$ be a general tuple of divisors, with
associated general complete intersection curve $C := D_1 \cap\cdots\cap D_{n-1}$. The
curve $C$ is then smooth and entirely contained in $X_{\reg}$. The restriction
$\sF|_C$ is locally free and nef, with $c_1(\sF|_C) = 0$. It follows that $\sF|_C$ is
semistable and therefore that $\sF$ is semistable with respect to $H$. But then,
$\Omega^{[1]}_X$ will likewise be semistable with respect to $H$.

\subsection*{Step 2: Local freeness of $\sF$}

Next, choose a resolution of singularities $\pi: \wtilde{X} \to X$ such that the
quotient sheaf
\[
\wtilde{\sF} := \factor{\pi^*(\sF)}{\tor}
\]
is locally free; such a resolution exists by \cite[Th.\,3.5]{Rossi68}. Recall
from Fact~\ref{fact:1} that since $\sF$ is nef, then so are $\pi^*(\sF)$, $\wtilde{\sF}$
and $\det \wtilde{\sF}$. There is more that we can say. Since the determinant of
$\sF$ is trivial by construction, $\det \sF \cong \sO_X$, we may write
\[
\det \wtilde{\sF} = \sO_{\wtilde X}\bigl(\ts\sum a_i\cdot E_i\bigr) \quad \text{where
$E_i$ are $\pi$-exceptional and }a_i \in \bZ.
\]
The divisor $\sum a_i\cdot E_i$ is nef. But then the Negativity Lemma,
\cite[Lem.\,3.39(1)]{KM98}, asserts that $a_i \leq 0$ for all $i$, which of course
means that all $a_i = 0$. It follows that $\det \wtilde{\sF} = \sO_{\wtilde X}$ and
thus that $\wtilde{\sF}$ is numerically flat. This has two consequences.
\begin{itemize}
\item First, it follows from the descent theorem for vector bundles on
resolutions of klt spaces, \cite[Th.\,1.2]{GKPT17}, that $\sF$ is locally free,
numerically flat, and that $\wtilde{\sF} = \pi^*(\sF)$.

\item Second, it follows from \cite[Prop.\,3.7]{GKP20a} that $\Omega^1_{X_{\reg}}$ is
projectively flat.
\end{itemize}

\subsection*{Step 3: Singularities of $X$}

If $\what{X} \to X$ is any maximally quasi-étale cover, then local freeness of $\sF$
immediately implies local freeness and nefness for the reflexive normalised
cotangent sheaf of $\what{X}$. The covering space $\what{X}$ therefore
reproduces the assumptions of Theorem~\ref{thm:6-1}, which allows us to assume
without loss of generality that $X$ is itself maximally quasi-étale. Together
with projective flatness of $\Omega^1_{X_{\reg}}$, this assumption allows to apply the
local description of projectively flat sheaves found in
\cite[Prop.\,3.11]{GKP20a}: every singular point $x \in X^{(\an)}$ admit a
neighbourhood $U$, open in the analytic topology, such that
\[
\Omega^{[1]}_U \simeq \sL_U^{\oplus n}
\]
with some Weil-divisorial sheaf $\sL_U$ on $U$. It follows from
Proposition~\ref{prop:icq} that $X$ has at worst isolated singularities.

\subsection*{Step 4: End of proof in case where $X$ is higher-dimensional}

If $n \geq 3$, consider the general complete intersection surface
$S := D_1 \cap\cdots\cap D_{n-2}$. Using that $X$ has isolated singularities, we find
that $S$ is smooth and contained in $X_{\reg}$, so that $\Omega^1_X|_S$ is locally
free, $H$-semistable and projectively flat. But by \cite[Prop.\,3.1.b on
p.\,42]{Kob87} this implies
\[
\frac{n-1}{2n}\cdot c_1\bigl(\Omega^1_X|_S\bigl)^2 = c_2\bigl(\Omega^1_X|_S\bigl)
\]
and therefore
\[
\frac{n-1}{2n}\cdot \what{c}_1(X)^2\cdot [H]^{n-2} = \what{c}_2(X)\cdot [H]^{n-2}.
\]
The claim thus follows from the semistability of $\Omega^1_X$ and
Theorem~\ref{thm:charQAbelian}.

\subsection*{Step 5: End of proof in case where $X$ is a surface}

It remains to consider the case where $\dim X = 2$, so that $X$ is a surface
with klt quotient singularities. We have seen above that
\[
\Sym^{[2]} \Omega^1_X = \sF \otimes \sO_X(K_X),
\]
where $\sF$ is a locally free numerically flat sheaf of rank three. The second
$\bQ$-Chern class is therefore computed as
\[
\what c_2 \bigl(\Sym^{[2]} \Omega^{[1]}_X \bigr) = \what{c}_2 \bigl(\sF \otimes \sO_X(K_X) \bigr) = \what{c}_2 \bigl(\sO_X(K_X)^{\oplus 3} \bigr) = 3\cdot [K_X]^2.
\]
On the other hand, the standard formula for the second Chern class of a symmetric product gives
\[
\what c_2 \bigl(\Sym^{[2]} \Omega^{[1]}_X \bigr) = 2\cdot [K_X]^2 + 4\cdot \what{c}_2(X).
\]
Comparing these two equations, we find $\frac{1}{4}\cdot [K_X]^2 = \what c_2(X)$.
Thanks to the semistability of $\Omega^{[1]}_X$, we may again apply
Theorem~\ref{thm:charQAbelian} and end the proof. \qed

\backmatter
\bibliographystyle{jepalpha+eid}
\bibliography{greb-et-al}
\end{document}

